% 第2章 流体运动学 —— 依据原书扫描页手工精修版
\chapter{流体运动学}

本章运用分析和几何描述方法来研究流体一般运动的性质，即流体运动学.至于产生流动的动力学原因将在后续各章讨论.

\section*{2.1  描述流体运动的两种方法}\addcontentsline{toc}{section}{2.1 描述流体运动的两种方法}

前面已经有了流体的连续介质模型和微团概念，流体的运动是无穷多微团的运动，要建立分析方法来精细地刻画流动过程中微团集合的运动状态.流体力学中有两种描述无穷多连续分布微团运动的方法.第一种方法是最常用的场描述方法，这种方法的基本思想是：在任意指定的时间逐点描述当地的运动特征量(如速度，加速度)及其他的物理量分布(如压力，密度等).这种方法称为欧拉描述法.第二种方法是跟踪质点的描述法，这种方法的基本思想是：从某个时刻开始跟踪每一个质点，记录这些质点的位置、速度、加速度和物理参数的变化，这种方法是离散质点运动描述方法在流体力学中的推广，称为拉格朗日描述法.下面建立这两种描述方法的分析公式，并加以比较.

\subsection*{1. 欧拉描述法}

欧拉描述法：在选定的时空坐标系$\{\boldsymbol{x},t\}$中考察流动过程中力学和其他物理参量的分布.时空坐标$\{\boldsymbol{x},t\}$是自变量，称作\textbf{欧拉变量}；当地的物理参量：如流动速度$\boldsymbol{U}$，温度$T$，压强$p$等表示为欧拉变量的函数，即
\begin{equation*}
\boldsymbol{U}=\boldsymbol{U}(\boldsymbol{x},t),\quad T=T(\boldsymbol{x},t),\quad p=p(\boldsymbol{x},t)
\tag{2.1}
\end{equation*}
欧拉方法是一种场的描述法，式(2.1)常称为流场中的速度分布，温度分布和压强分布等，或简称速度场、温度场和压强场.

\subsection*{2. 拉格朗日描述法}

拉格朗日描述法是跟踪质点来描述它们的力学和其他物理状态.实现这种方法的关键是建立识别质点的方法.最方便的方法是用每个质点在初始时刻的坐标作为它们的“标记”，然后跟踪每个质点，在它们的运动轨迹上考察它们的物理状态.连续介质可分割成无限多连绵一片的质点，因此连续介质质点的初始时刻坐标$\boldsymbol{A}(a_1,a_2,a_3)$在考察的区域内是连续分布的.质点的初始时刻坐标$\boldsymbol{A}$和时间变量$t$是拉格朗日法的自变量，称为\textbf{拉格朗日变量}.流体质点的位移$\boldsymbol{x}$，温度$T$和压强$p$等是拉格朗日变量的函数，即
\begin{equation*}
\boldsymbol{x}=\boldsymbol{x}(\boldsymbol{A},t),\quad T=T(\boldsymbol{A},t),\quad p=p(\boldsymbol{A},t)
\tag{2.2}
\end{equation*}
拉格朗日描述法中的位移函数$\boldsymbol{x}=\boldsymbol{x}(\boldsymbol{A},t)$，也就是质点的轨迹，有以下两个基本性质.

(1)在初始时刻$t=0$时，$\boldsymbol{x}=\boldsymbol{A}$，即
\begin{equation*}
\boldsymbol{x}(\boldsymbol{A},0)=\boldsymbol{A}
\tag{2.3}
\end{equation*}

(2)在任何时刻，质点位置变量$\boldsymbol{x}$与该质点初始时刻的位置变量$\boldsymbol{A}$是一一对应的连续函数.

性质(1)是拉格朗日变量$\boldsymbol{A}$的定义，性质(2)是连续介质的属性.由于变量$\boldsymbol{x}$和变量$\boldsymbol{A}$是一一对应的连续函数，根据数学分析中的隐函数定理，$\boldsymbol{x}=\boldsymbol{x}(\boldsymbol{A},t)$必存在反演
\begin{equation*}
\boldsymbol{A}=\boldsymbol{A}(\boldsymbol{x},t)
\tag{2.4}
\end{equation*}
反演式$\boldsymbol{A}=\boldsymbol{A}(\boldsymbol{x},t)$的意义是：在$t$时刻位于$\boldsymbol{x}$的质点可由式(2.4)追溯到它的初始位置$\boldsymbol{A}$.就是说式(2.4)可将描述质点的欧拉变量$\{\boldsymbol{x},t\}$转换成拉格朗日变量$\{\boldsymbol{A},t\}$，而式(2.2)中的第一式是将描述质点的拉格朗日变量转换成的欧拉变量.图2.1示意欧拉变量$\{\boldsymbol{x},t\}$和拉格朗日变量$\{\boldsymbol{A},t\}$的转换.

\subsection*{3. 欧拉描述式和拉格朗日表达式的互换}

流动过程的物理量演化既可用欧拉方法表达也可以用拉格朗日方法表达，它们之间可以互相转换.

(1)拉格朗日描述式变换成欧拉描述式(L-E变换)

L-E变换就是把物理量在时空中演化的拉格朗日描述式替换成欧拉描述式.设已知拉格朗日的位移表达式：
\[
\boldsymbol{x}=\boldsymbol{x}(\boldsymbol{A},t)
\]
和压强、密度表达式
\[
p=p(\boldsymbol{A},t),\qquad \rho=\rho(\boldsymbol{A},t)
\]
要求把$p$，$\rho$的拉格朗日表达式转换成欧拉表达式.拉格朗日位移$\boldsymbol{x}=\boldsymbol{x}(\boldsymbol{A},t)$有唯一的反演式：$\boldsymbol{A}=\boldsymbol{A}(\boldsymbol{x},t)$(式(2.4))，将它代入压强或密度的拉格朗日表达式，就可以得到相应变量的欧拉表达式，例如，压强的拉格朗日表达式为$p=p_L(\boldsymbol{A},t)$，将式(2.4)代入，
\[
p=p_L(\boldsymbol{A},t)=p_L[\boldsymbol{A}(\boldsymbol{x},t),t]=p_E(\boldsymbol{x},t)
\]
得到压强的欧拉表达式.这里，为了明确起见，把拉格朗日描述式记以下标L，欧拉描述式记以下标E.密度及其他物理参量的转换方法与此相同.

由拉格朗日位移表达式求得欧拉的速度场表达式，需要先求出速度的拉格朗日表达式，然后再用上面的方法做变量替换.质点速度是质点的位移对时间的导数，对于指定质点，它的拉格朗日坐标$\boldsymbol{A}$是常量，即质点标记不变，因此拉格朗日描述法中质点速度等于位移函数对时间求一阶偏导数，即
\begin{equation*}
\boldsymbol{U}=\Big[\frac{\partial\boldsymbol{x}(\boldsymbol{A},t)}{\partial t}\Big]_{\boldsymbol{A}}=\boldsymbol{U}_L(\boldsymbol{A},t)
\tag{2.5}
\end{equation*}
用式(2.4)$\boldsymbol{A}=\boldsymbol{A}(\boldsymbol{x},t)$替换式(2.5)中$\boldsymbol{A}$，就得到速度场的欧拉描述式：
\[
\boldsymbol{U}=\boldsymbol{U}_L(\boldsymbol{A},t)=\boldsymbol{U}_L[\boldsymbol{A}(\boldsymbol{x},t),t]=\boldsymbol{U}_E(\boldsymbol{x},t)
\]

\noindent\textbf{例2.1}\quad 给定拉格朗日位移描述式：
\[
x_1=a_1\exp\Big(-\frac{2t}{k}\Big),\quad x_2=a_2\exp\Big(\frac{t}{k}\Big),\quad x_3=a_3\exp\Big(\frac{t}{k}\Big)
\]
求欧拉速度场.$k$是常数，$\{a_i,t\}$是拉格朗日变量.

\noindent\textbf{解}\quad 先求速度的拉格朗日表达式
\begin{gather*}
U_1=\Big(\frac{\partial x_1}{\partial t}\Big)_{\!A}=-\frac{2a_1}{k}\exp\Big(-\frac{2t}{k}\Big)\\[2pt]
U_2=\Big(\frac{\partial x_2}{\partial t}\Big)_{\!A}=\frac{a_2}{k}\exp\Big(\frac{t}{k}\Big)\\[2pt]
U_3=\Big(\frac{\partial x_3}{\partial t}\Big)_{\!A}=\frac{a_3}{k}\exp\Big(\frac{t}{k}\Big)
\end{gather*}
然后由位移表达求反演，
\[
a_1=x_1\exp\Big(\frac{2t}{k}\Big),\quad a_2=x_2\exp\Big(-\frac{t}{k}\Big),\quad a_3=x_3\exp\Big(-\frac{t}{k}\Big)
\]
将$a_1,a_2,a_3$代入速度式，得速度场的欧拉表达式.
\[
U_1=-\frac{2}{k}x_1,\qquad U_2=\frac{1}{k}x_2,\qquad U_3=\frac{1}{k}x_3
\]

(2)欧拉描述式转换成拉格朗日描述式(E-L变换)

E-L变换是把流场表达式中欧拉变量$(\boldsymbol{x},t)$用拉格朗日变量$(\boldsymbol{A},t)$替换，也就是首先需要求得质点的位移函数$\boldsymbol{x}=\boldsymbol{x}(\boldsymbol{A},t)$.设已知欧拉速度场表达式：$\boldsymbol{U}=\boldsymbol{U}_E(\boldsymbol{x},t)$，我们要求它的拉格朗日表达式.

由于速度是质点位移向量对时间的一阶偏导数，在拉格朗日描述法中：$\boldsymbol{U}=\Big(\dfrac{\partial\boldsymbol{x}}{\partial t}\Big)_{\!A}$，因此首先要从下式求出质点的位移函数：
\begin{equation*}
\boldsymbol{U}=\Big(\frac{\partial\boldsymbol{x}}{\partial t}\Big)_{\!A}=\boldsymbol{U}_E(\boldsymbol{x},t)
\tag{2.6}
\end{equation*}
式(2.6)是关于质点位移$\boldsymbol{x}(t)$的一阶常微分方程组，它的初始条件为
\begin{equation*}
t=0:\ \boldsymbol{x}=\boldsymbol{A}
\tag{2.7}
\end{equation*}
积分式(2.6)，得到质点位移的表达式：
\[
\boldsymbol{x}=\boldsymbol{x}(\boldsymbol{A},t)
\]
代入速度表达式，得到速度的拉格朗日描述式：
\[
\boldsymbol{U}=\boldsymbol{U}_E[\boldsymbol{x}(\boldsymbol{A},t),t]=\boldsymbol{U}_L(\boldsymbol{A},t)
\]

\noindent\textbf{例2.2}\quad 已知欧拉速度场：$U_1=x_1+t$，$U_2=x_2+t$，$U_3=0$，求质点位移和速度的拉格朗日表达式.

\noindent\textbf{解}\quad 解常微分方程
\[
U_1=\frac{\partial x_1}{\partial t}=x_1+t,\qquad U_2=\frac{\partial x_2}{\partial t}=x_2+t,\qquad U_3=\frac{\partial x_3}{\partial t}=0
\]
初始条件为：$t=0$，$x_1=a_1$，$x_2=a_2$，$x_3=a_3$.该微分方程的一般解为
\[
x_1=c_1\exp(t)-t-1,\qquad x_2=c_2\exp(t)-t-1,\qquad x_3=c_3
\]
将初始条件代入后可得积分常数$c_1=a_1+1$，$c_2=a_2+1$，$c_3=a_3$，最后拉格朗日位移表达式为
\[
x_1=(a_1+1)\exp(t)-t-1,\qquad x_2=(a_2+1)\exp(t)-t-1,\qquad x_3=a_3
\]
将位移表达式代入速度场公式，得速度的拉格朗日描述式如下：
\[
U_1=(a_1+1)\exp(t)-1,\qquad U_2=(a_2+1)\exp(t)-1,\qquad U_3=0
\]
\section*{2.2  流场的几何描述}\addcontentsline{toc}{section}{2.2 流场的几何描述}

为了直观和形象地了解流动状态，可以用几何方法来描述流场，例如图2.2是气流绕机翼剖面(又称翼型)的流动.本来气体是透明的，无法辨认它的运动状态，在气流的上游注入几股有色的烟丝，烟丝被气流携带，它们的流速等于当时当地气体速度，如果用曝光很快的相机摄下烟丝的图像，它们就反映某一瞬间翼型周围气流运动的方向，这就是一种显示气流运动的方法，烟丝曲线称为气流的\textbf{流线}.

下面给出流线的数学定义，并进一步分析它的性质.

\subsection*{1. 流线、流面和流管}

\noindent\textbf{定义2.1}\quad 流线是速度场的向量线.它是某一固定时刻的空间曲线，该曲线上任意一点的切向量与当地的速度向量重合.

根据定义，流线是在给定时间考察流动在空间的图像，它是一种欧拉速度场的描述方法.

设流线的参数方程为
\[
\boldsymbol{x}=\boldsymbol{r}(s)
\]
则流线任意点处切向量为$\mathrm{d}\boldsymbol{r}$，流场当地速度为$\boldsymbol{U}(\boldsymbol{x},t_0)$(欧拉表达式)，根据流线的定义，我们要求$\mathrm{d}\boldsymbol{r}/\!/\boldsymbol{U}(\boldsymbol{x},t_0)$，应用向量代数运算公式，流线上应满足
\begin{equation*}
\mathrm{d}\boldsymbol{r}\times\boldsymbol{U}=\boldsymbol{0}
\tag{2.8}
\end{equation*}
展开式(2.8)，并简化后可得
\begin{equation*}
\frac{\mathrm{d}x_1}{U_1(x_1,x_2,x_3,t_0)}=\frac{\mathrm{d}x_2}{U_2(x_1,x_2,x_3,t_0)}=\frac{\mathrm{d}x_3}{U_3(x_1,x_2,x_3,t_0)}
\tag{2.9}
\end{equation*}
给定流线的起始点：即$s=0$时$x_1=x_0$，$x_2=y_0$，$x_3=z_0$，积分式(2.9)就得到某时刻通过$(x_0,y_0,z_0)$的流线.必须强调指出，流线是在固定时刻对流场的描述，流线方程中时间$t_0$是常参量.也就是说，当需要考察某一时刻的流场时，就把该时刻$t_0$代入速度场表达式.

\noindent\textbf{例2.3}\quad 给定速度场：
\[
u=x+t,\qquad v=-y+t,\qquad w=0
\]
求$t=0$时，通过$M(-1,-1,0)$点的流线.

\noindent\textbf{解}\quad 流线的微分方程为
\[
\frac{\mathrm{d}x}{u}=\frac{\mathrm{d}y}{v}=\frac{\mathrm{d}z}{w}
\]
因$w=0$，故$\mathrm{d}z=0$，它的解是$z=$常数，将初始值$z_0=0$代入，得：$z=0$；这表明流线位于$z=0$的坐标面上.由于$u$，$v$只和$x$，$y$有关，因此流线方程可直接积分，将$t=0$代入速度表达式，得流线方程为
\[
\frac{\mathrm{d}x}{x}=\frac{\mathrm{d}y}{-y}
\]
积分后得
\[
xy=\mathrm{const.}
\]
将起始点$M(-1,-1,0)$代入积分式得：const.$=1$，该流线方程为
\[
xy=1,\qquad z=0
\]
本例的结果是流场在$t=0$时通过$M(-1,-1,0)$的流线是平面双曲线.

为直观了解某一流动状态，我们不仅要看过某点的一条流线，还需要了解同一时刻通过若干点的流线(如图2.2所示的烟丝)，通过若干点的流线族称为\textbf{流谱}.

\noindent\textbf{例2.4}\quad 试求：任意时刻例2.3给定的速度场的流线谱.

\noindent\textbf{解}\quad 求任意时刻的流线谱时，速度表达式中$t$作为参量，它在式(2.9)的积分中保持常数.因$w=0$，仍积分得$z=$常数，由
\[
\frac{\mathrm{d}x}{x+t}=\frac{\mathrm{d}y}{-y+t}
\]
积分函数为
\[
(x+t)(-y+t)=\mathrm{const.}
\]
通过以上积分式，可以考察不同时刻的流线谱.

(1)$t=0$时的流线谱：将$t=0$代入积分式，得$xy=\mathrm{const.}$，它是一族双曲线.其中通过$M(-1,-1,0)$点的流线是$xy=1$，它和例2.3结果相同.给定不同的常数值，得到一族双曲线.令const.$=0$，可得$t=0$时另外两条流线：$x=0\,(-\infty<y<\infty)$和$y=0\,(-\infty<x<\infty)$，它们是退化的双曲线.$t=0$的流线谱示于图2.3(a).$t=0$时的速度场是：$u=x$，$v=-y$，$w=0$.因此$x>0$(第一、第四象限)处的流场速度$u$朝$x$轴正方向；$x<0$(第二、第三象限)处的流场速度$u$朝$x$轴的负方向，而$y>0$(第一、第二象限)处的速度场$v$朝$y$轴负方向；$y<0$(第三、第四象限)处速度场$v$朝$y$轴正方向.图2.3(a)中给出了流线谱上的速度方向.

(2)$t=1$时的流线谱：将$t=1$代入积分式，得$(x+1)(y-1)=\mathrm{const.}$，它仍是一族双曲线，但是双曲线的中心移到$(-1,1)$.流线谱上仍标明速度方向.对于例题给定的流场，不同时刻流线谱只是发生平移.但是观察通过$M(-1,-1,0)$点(图2.3(a)、图2.3(b)上“$+$”号)的流线时，可以发现：在$t=0$时，通过该点的流线是双曲线；而$t=1$时，通过该点的流线是向上的直线.

从上面的实例中可以发现流线有以下性质.

第一，不同时刻通过同一点的流线可以不重合.一般情况下，如果速度场与时间有关，则不同时刻、通过同一点的速度大小和方向都可能是不同的.于是，从不同时刻、同一点出发的流线初始方向不同，这些流线当然不会重合.例如例2.2中$t=0$时通过$(-1,-1,0)$的流线是双曲线；而$t=1$时通过该点的流线是平行于$y$轴的直线.根据以上分析，还可以推断：如果欧拉速度场和时间无关，那么任何时刻通过同一点的流线都将是相同的.

第二，流线是有走向的几何线，流线的走向由速度场确定，通常在流线上任取一点，用该点的速度方向标明流线的走向，如图2.3所示.

第三，一般情况下，同一时刻通过一点只有一根流线.因为同一时刻，在空间一点上只有一个速度.换句话说，一般情况下，同一时刻流场中的流线不能相交.但是，在理论上有两种例外，\textcircled{1}流场中速度等于零的点，因为速度等于零的点其速度方向是任意的，不同方向的流线都可以通过该点；\textcircled{2}流场中速度为无限大的点，因为这种点上流体的速度方向也是不定的.例如，例2.4中$t=0$时刻有通过原点的两条相交直线；$t=1$时有通过$(-1,1)$点的相交流线，但是相交点的速度等于零.

如果给定另一种速度场：
\[
u=\frac{x}{\sqrt{x^{2}+y^{2}}},\qquad v=\frac{y}{\sqrt{x^{2}+y^{2}}}
\]
则不难求出它的流线谱，如图2.4所示，这时过原点$(0,0,0)$的流线相交，但是相交点的速度是无穷大.数学分析上允许流场中有无限大的速度，但是物理上是不存在无限大速度的.因此实际流场中发现流线相交时，相交点的速度一定等于零.

由流线概念可以扩充到流面和流管的概念.通过不在同一流线上的若干点作流线，构成流面或流管.

\noindent\textbf{定义2.2}\quad 某一时刻通过给定曲线(该曲线不是流线)上每一点作流线所构成的曲面称为流面.

简言之，流面是由流线组成的空间曲面.若给定的空间曲线为封闭曲线，则构成的流面是管状曲面，这种管状曲面的内域称为\textbf{流管}(图2.5).

流面由流线组成，因此流面上任意一点的法向量和当地的流场速度向量垂直，即在流面上有
\begin{equation*}
(\boldsymbol{U}\cdot\boldsymbol{n})_{\text{流面}}=0
\tag{2.10}
\end{equation*}
式中$\boldsymbol{n}$为流面法线方向.式(2.10)说明，流体不能穿过流面，也不能通过流管表面进入流管，所以流面如同刚性壁面一样，流体不能穿透它.

\subsection*{2. 迹线}

\noindent\textbf{定义2.3}\quad 流体质点的运动轨迹称为迹线.

迹线是流体质点运动的几何描述，在拉格朗日表达式中位移函数就是质点群的轨迹族：
\[
\boldsymbol{x}=\boldsymbol{x}(a,b,c,t)
\]
给定拉格朗日坐标$\boldsymbol{A}(a_1,a_2,a_3)$值就得到该质点的轨迹.

欧拉描述方法中流体质点轨迹需由速度场积分求出，如给定欧拉速度场$\boldsymbol{U}(\boldsymbol{x},t)$，则轨迹方程为
\begin{equation*}
\Big(\frac{\partial\boldsymbol{x}}{\partial t}\Big)_{\!A}=\boldsymbol{U}(\boldsymbol{x},t)
\tag{2.11a}
\end{equation*}
或
\begin{equation*}
\begin{cases}
\Big(\dfrac{\partial x_1}{\partial t}\Big)_{\!A}=U_1(x_1,x_2,x_3,t)\\[6pt]
\Big(\dfrac{\partial x_2}{\partial t}\Big)_{\!A}=U_2(x_1,x_2,x_3,t)\\[6pt]
\Big(\dfrac{\partial x_3}{\partial t}\Big)_{\!A}=U_3(x_1,x_2,x_3,t)
\end{cases}
\tag{2.11b}
\end{equation*}
在迹线方程中，$\boldsymbol{x}$或$(x_1,x_2,x_3)$是质点的坐标，它是时间$t$的函数.给定起始时刻$t=0$时质点的坐标$(a_1,a_2,a_3)$，积分式(2.11)就得到该质点的轨迹.

\noindent\textbf{例2.5}\quad 给定欧拉速度场：$U_1=x_1+t$，$U_2=-x_2+t$，$U_3=0$，求$t=0$时位于$M(-1,-1,0)$处的质点运动轨迹.

\noindent\textbf{解}\quad 写出迹线方程(略去式(2.11b)下标$A$)
\[
U_1=\frac{\partial x_1}{\partial t}=x_1+t,\qquad U_2=\frac{\partial x_2}{\partial t}=-x_2+t,\qquad U_3=\frac{\partial x_3}{\partial t}=0
\]
注意迹线方程中$t$是积分的自变量，$x_i$是$t$的函数，积分该方程后得
\[
x_1=C_1\exp(t)-t-1,\qquad x_2=C_2\exp(-t)+t-1,\qquad x_3=C_3
\]
将初始坐标代入后得：$C_1=C_2=C_3=0$，因此
\[
x_1=-t-1,\qquad x_2=t-1,\qquad x_3=0
\]
消去$x_1$，$x_2$方程中的$t$，得
\[
x_1+x_2=-2
\]
将该轨迹绘于图2.6.注意例2.5的速度场和起始点与例2.3相同，在同一图上绘出图2.3中$t=0$时过同一点的流线以作比较.我们可以看到：通过一点的流线与迹线在起始点相切，而后开始分叉，就是说例2.3给定的速度场中通过同一点的流线和轨迹不重合.这种情况在一般流场中是普遍的，本书将在下面详细讨论.

综合以上分析，流线和迹线是描述流场的不同几何特性，它们的最基本差别是：

\textbf{迹线是同一质点在不同时刻的位移曲线；}

\textbf{流线是同一时刻、不同质点连接起来的速度场向量线.}

简言之，迹线是描述指定质点的运动过程，流线是描述给定瞬间的速度场状态.

\subsection*{3. 定常流场和非定常流场}

\noindent\textbf{定义2.4}\quad 与时间无关的欧拉速度场称为定常流场，反之则称为非定常流场.

按照定义，定常速度场的表达式中不含欧拉变量$t$，只有位置变量$\boldsymbol{x}$，即
\[
\boldsymbol{U}=\boldsymbol{U}(\boldsymbol{x})
\]
定常流场与非定常流场是流体运动学中经常遇到的基本概念.最简单的定常流动例子是以恒定的压强差在管道中输送液体.如输送液体时压强差发生波动，则流动就不定常.一般来说非定常流动较定常流动复杂.流场是否定常还与观察者所在的坐标系有关，例如，在匀速水平飞行的飞弹中观察飞弹周围的气流，这时气流绕过飞弹是定常流场(图2.7(a))，而从地面静止的观察者看到的飞弹周围流场是非定常流场(图2.7(b))，他看到飞弹头部不断地把空气排开，而尾部有空气“吸入”，流线谱随着飞弹向前推进，是不定常的.

定常流场的流线和迹线有以下性质.

\textcircled{1} 在定常流场中通过同一点的流线不随时间变化；

\textcircled{2} 任意时刻通过同一空间点的迹线和流线重合.

可以由流线和迹线方程来证明上述性质.对于定常流动，速度场不含变量$t$，故流线方程和迹线方程及相应起始条件分别为
\[
\text{流线：}\ \frac{\mathrm{d}\boldsymbol{x}}{\mathrm{d}s}=\boldsymbol{U}(\boldsymbol{x})\,;\quad s=0\,,\ \boldsymbol{x}=\boldsymbol{A}
\]
\[
\text{迹线：}\ \Big(\frac{\partial\boldsymbol{x}}{\partial t}\Big)_{\!A}=\boldsymbol{U}(\boldsymbol{x})\,;\quad t=0\,,\ \boldsymbol{x}=\boldsymbol{A}
\]
流线方程右端不含时间，因此流线方程的积分曲线$\boldsymbol{x}=\boldsymbol{r}(s,\boldsymbol{A})$与时间无关，这就证实了定常流的第一个运动学性质.对于迹线来说，它的方程右端也不含时间$t$，因此迹线上任意一点的切向量与当地流线的切向量处处重合，从而由同一点出发的迹线积分曲线也必定和流线重合.这就证明了定常流动的第二个运动学性质.从数学上来看，定常流动中流线和迹线有相同的常微分方程和初始条件(只是自变量用不同符号)，积分曲线当然重合.

一般情况下，非定常流场中流线谱随时间变化，而且通过同一点的流线和迹线不再重合，如例2.5.但是，有个别的非定常运动，它们的流场速度方向不随时间变化，只有速度大小随时间变化，这时流线和迹线仍然可能重合.

\subsection*{4. 流体线与流体面及其保持性}

在流体力学中常常还需要分析由相同质点组成的几何面，例如：液滴表面、气泡表面以及水波与空气的界面等.当流体运动时，这种界面也随着运动和变形，但界面始终是由给定流体质点组成.

\noindent\textbf{定义2.5}\quad 由同一组连续质点组成的几何曲线或曲面分别称为流体质线或流体质面，简称流体线或流体面.

由于流体的连续性质，流体线或流体面在运动过程中可以变形，但不能断裂，也就是质点间的相邻关系不能改变，这一性质称作流体面的\textbf{保持性}.下面将导出流体面保持性的数学表达式.设有一流体面，它的几何方程为
\begin{equation*}
F(\boldsymbol{x},t)=0
\tag{2.12}
\end{equation*}
我们将证明，流体面方程必须满足以下关系式：
\begin{equation*}
\frac{\partial F}{\partial t}+\frac{\partial F}{\partial x_1}\frac{\delta x_1}{\delta t}+\frac{\partial F}{\partial x_2}\frac{\delta x_2}{\delta t}+\frac{\partial F}{\partial x_3}\frac{\delta x_3}{\delta t}=\frac{\partial F}{\partial t}+\boldsymbol{U}\cdot\nabla F=0
\tag{2.13}
\end{equation*}
式中$\boldsymbol{U}$是流体面上的流体质点速度.式(2.13)证明如下.设流体面上的一质点在$\delta t$时间内的位移为$\delta\boldsymbol{x}$，因为该质点在流体面上，所以到达新的时空坐标$(t+\delta t,\boldsymbol{x}+\delta\boldsymbol{x})$时，时空坐标仍然满足式(2.12)，即
\[
F(\boldsymbol{x}+\delta\boldsymbol{x},t+\delta t)=0
\]
利用泰勒展开式
\begin{align*}
F(\boldsymbol{x}+\delta\boldsymbol{x},t+\delta t)&=F(\boldsymbol{x},t)+\frac{\partial F}{\partial t}\delta t+\frac{\partial F}{\partial x_1}\delta x_1+\frac{\partial F}{\partial x_2}\delta x_2\\
&\quad+\frac{\partial F}{\partial x_3}\delta x_3+O(\delta t^{2},|\delta\boldsymbol{x}|^{2},\delta t|\delta\boldsymbol{x}|)
\end{align*}
式中$F(\boldsymbol{x}+\delta\boldsymbol{x},t+\delta t)=F(\boldsymbol{x},t)=0$.将上式除以$\delta t$并令$\delta t\to 0$，$|\delta\boldsymbol{x}|\to 0$，得
\[
\frac{\partial F}{\partial t}+\frac{\partial F}{\partial x_1}\frac{\delta x_1}{\delta t}+\frac{\partial F}{\partial x_2}\frac{\delta x_2}{\delta t}+\frac{\partial F}{\partial x_3}\frac{\delta x_3}{\delta t}=0
\]
因为$(\delta x_1,\delta x_2,\delta x_3)$是质点的位移，故
\[
\frac{\delta x_1}{\delta t}=U_1,\qquad \frac{\delta x_2}{\delta t}=U_2,\qquad \frac{\delta x_3}{\delta t}=U_3
\]
它们是质点的速度，将它代入前式，得式(2.13)，它是流体面保持性的数学表达式.
\section*{2.3  质点的加速度公式和质点导数}\addcontentsline{toc}{section}{2.3 质点的加速度公式和质点导数}

\subsection*{1. 质点加速度公式}

质点加速度是质点速度向量随时间的变化率.在拉格朗日描述式中，位移函数$\boldsymbol{x}(\boldsymbol{A},t)$本来就追踪质点，并以$\boldsymbol{A}$识别某一质点，因此位移函数对时间$t$求二次偏导数，就得到质点$\boldsymbol{A}$的加速度
\begin{equation*}
\boldsymbol{a}=\Big(\frac{\partial^{2}\boldsymbol{x}}{\partial t^{2}}\Big)_{\!A}
\tag{2.14}
\end{equation*}
欧拉描述法给定速度场$\boldsymbol{U}(\boldsymbol{x},t)$，而不追踪质点，要由速度场计算$(\boldsymbol{x},t)$处的质点加速度时必须求出该质点在$\delta t$时间内的速度增量，然后求极值，即
\[
\boldsymbol{a}=\lim_{\substack{\delta\boldsymbol{x}\to 0\\ \delta t\to 0}}\bigg[\frac{\boldsymbol{U}(\boldsymbol{x}+\delta\boldsymbol{x},t+\delta t)-\boldsymbol{U}(\boldsymbol{x},t)}{\delta t}\bigg]
\]
式中$\delta\boldsymbol{x}$是质点在$\delta t$时间内的位移，将加速度计算公式的分子作泰勒展开，
\begin{align*}
\boldsymbol{U}(\boldsymbol{x}+\delta\boldsymbol{x},t+\delta t)-\boldsymbol{U}(\boldsymbol{x},t)&=\Big(\frac{\partial\boldsymbol{U}}{\partial t}\Big)_{\!x}\delta t+\Big(\frac{\partial\boldsymbol{U}}{\partial x_1}\Big)_{\!x}\delta x_1+\Big(\frac{\partial\boldsymbol{U}}{\partial x_2}\Big)_{\!x}\delta x_2\\
&\quad+\Big(\frac{\partial\boldsymbol{U}}{\partial x_3}\Big)_{\!x}\delta x_3+O(\delta t^{2},|\delta\boldsymbol{x}|^{2},\delta t|\delta\boldsymbol{x}|)
\end{align*}
将它代入加速度公式，并略去高阶小量后得
\[
\boldsymbol{a}=\frac{\partial\boldsymbol{U}}{\partial t}+\frac{\partial\boldsymbol{U}}{\partial x_1}\frac{\delta x_1}{\delta t}+\frac{\partial\boldsymbol{U}}{\partial x_2}\frac{\delta x_2}{\delta t}+\frac{\partial\boldsymbol{U}}{\partial x_3}\frac{\delta x_3}{\delta t}
\]
注意到$\delta\boldsymbol{x}$是质点位移，因而
\[
\lim_{\delta t\to 0}\frac{\delta x_1}{\delta t}=U_1,\qquad \lim_{\delta t\to 0}\frac{\delta x_2}{\delta t}=U_2,\qquad \lim_{\delta t\to 0}\frac{\delta x_3}{\delta t}=U_3
\]
于是，加速度公式为
\begin{equation*}
\boldsymbol{a}=\frac{\partial\boldsymbol{U}}{\partial t}+U_1\frac{\partial\boldsymbol{U}}{\partial x_1}+U_2\frac{\partial\boldsymbol{U}}{\partial x_2}+U_3\frac{\partial\boldsymbol{U}}{\partial x_3}
\tag{2.15}
\end{equation*}
利用向量算符$\nabla=\boldsymbol{e}_1\dfrac{\partial}{\partial x_1}+\boldsymbol{e}_2\dfrac{\partial}{\partial x_2}+\boldsymbol{e}_3\dfrac{\partial}{\partial x_3}$，可得$\boldsymbol{U}\cdot\nabla=U_1\dfrac{\partial}{\partial x_1}+U_2\dfrac{\partial}{\partial x_2}+U_3\dfrac{\partial}{\partial x_3}$，式(2.15)还可写作
\begin{equation*}
\boldsymbol{a}=\frac{\partial\boldsymbol{U}}{\partial t}+(\boldsymbol{U}\cdot\nabla)\boldsymbol{U}
\tag{2.16}
\end{equation*}
在直角坐标中加速度的分量式(用习惯的$x,y,z$符号)为
\begin{gather*}
a_x=\frac{\partial u}{\partial t}+u\frac{\partial u}{\partial x}+v\frac{\partial u}{\partial y}+w\frac{\partial u}{\partial z}\tag{2.17a}\\[2pt]
a_y=\frac{\partial v}{\partial t}+u\frac{\partial v}{\partial x}+v\frac{\partial v}{\partial y}+w\frac{\partial v}{\partial z}\tag{2.17b}\\[2pt]
a_z=\frac{\partial w}{\partial t}+u\frac{\partial w}{\partial x}+v\frac{\partial w}{\partial y}+w\frac{\partial w}{\partial z}\tag{2.17c}
\end{gather*}
从加速度公式可以看出，用欧拉速度场计算质点运动速度随时间的变化率时，质点加速度由两部分组成：

(1)\textbf{局部导数}$\dfrac{\partial\boldsymbol{U}}{\partial t}$，它是给定空间点上的速度随时间的变化率，这一部分是因流场的非定常性产生的；

(2)\textbf{对流导数或迁移导数}$(\boldsymbol{U}\cdot\nabla)\boldsymbol{U}$.因为$\boldsymbol{U}\cdot\nabla=U\,\dfrac{\partial}{\partial s}$($s$是速度方向)，它是沿速度方向的导数，或给定瞬时的速度场流线方向的导数，这一部分是因速度场不均匀性产生的.

于是，可得如下结论：用欧拉法给出的速度场求加速度时，

\textbf{质点加速度$=$速度的局部导数$+$速度的迁移导数}

\subsection*{2. 质点导数}

将推导加速度公式的方法推广到质点上任意物理量增长率的计算，就有质点导数的概念.

\noindent\textbf{定义2.6}\quad 质点携带的物理量随时间的变化率称为质点导数，并用$\dfrac{D}{Dt}$表示.

不难得到，在欧拉描述法中任意物理量$Q$的质点导数等于
\begin{equation*}
\frac{DQ}{Dt}=\frac{\partial Q}{\partial t}+\boldsymbol{U}\cdot\nabla Q
\tag{2.18}
\end{equation*}
上式证明如下：任意物理量的欧拉表达式为$Q=Q(\boldsymbol{x},t)$，$Q$可以是标量、向量或张量，它的质点导数等于：
\[
\frac{DQ}{Dt}=\lim_{\substack{\delta\boldsymbol{x}\to 0\\ \delta t\to 0}}\bigg[\frac{Q(\boldsymbol{x}+\delta\boldsymbol{x},t+\delta t)-Q(\boldsymbol{x},t)}{\delta t}\bigg]
\]
利用泰勒展开
\[
Q(\boldsymbol{x}+\delta\boldsymbol{x},t+\delta t)-Q(\boldsymbol{x},t)=\Big(\frac{\partial Q}{\partial t}\Big)_{\!x}\delta t+\delta\boldsymbol{x}\cdot\nabla Q+O(\delta t^{2},|\delta\boldsymbol{x}|^{2},\delta t|\delta\boldsymbol{x}|)
\]
代入前面公式就得式(2.18).质点导数公式对任意物理量都成立，因此可以将质点导数的运算用以下算符表示
\[
\frac{D}{Dt}=\frac{\partial}{\partial t}+\boldsymbol{U}\cdot\nabla=\frac{\partial}{\partial t}+U_1\frac{\partial}{\partial x_1}+U_2\frac{\partial}{\partial x_2}+U_3\frac{\partial}{\partial x_3}
\]
在欧拉描述法中的质点导数可用文字表述如下：

\textbf{物理量的质点导数$=$物理量的局部导数$+$物理量的对流导数.}

在欧拉描述法中必须将质点导数$\dfrac{D}{Dt}$和欧拉场量的局部导数$\dfrac{\partial}{\partial t}$明确区别开，虽然两者都是对时间导数的形式.局部导数是物理量在欧拉变量$\boldsymbol{x}=(x_1,x_2,x_3)$不变的条件下的时间增长率；而质点导数是物理量在质点运动轨迹上的时间增长率.

定常场中，物理量的欧拉表达式不含时间变量，因而局部导数等于零，但是质点导数可以不等于零，它等于对流导数.就是说：\textbf{定常流场中物理量的质点导数等于物理量的对流导数.}

用欧拉法描述的具体流动现象中，局部导数是物理量在当地的变化率，例如某地气象温度的变化率，或河道某处的水位增长率等都属于局部导数.如果要知道流动过程中物理量的质点导数，例如质点的密度或温度的变化率等，除了有该物理量的欧拉表达式外，还必须知道速度场的欧拉表达式，分别计算出该物理量的局部导数和对流导数，然后用式(2.18)计算质点导数.

\noindent\textbf{例2.6}\quad 由气象观测站测得的大气温度和速度分布如下：
\[
\boldsymbol{U}=U(y)\boldsymbol{e}_x,\qquad T=T_0(x)+\alpha\exp(-\gamma t^{2})
\]
求在$(x_0,y_0,z_0;t_0)$处温度的质点导数.

\noindent\textbf{解}\quad 由质点导数公式(2.18)得
\[
\frac{DT}{Dt}=\frac{\partial T}{\partial t}+u\frac{\partial T}{\partial x}+v\frac{\partial T}{\partial y}+w\frac{\partial T}{\partial z}
\]
其中各导数等于
\begin{gather*}
\frac{\partial T}{\partial t}=-2\alpha\gamma t\exp(-\gamma t^{2}),\qquad \frac{\partial T}{\partial y}=\frac{\partial T}{\partial z}=0,\qquad \frac{\partial T}{\partial x}=T_0'(x)\\[4pt]
\frac{DT}{Dt}=-2\alpha\gamma t\exp(-\gamma t^{2})+U(y)\,T_0'(x)
\end{gather*}
将$(x_0,y_0,z_0;t_0)$代入公式右边的各项，就可得该点温度的质点导数：
\[
\frac{DT}{Dt}(x_0,y_0,z_0;t_0)=-2\alpha\gamma t_0\exp(-\gamma t_0^{2})+U(y_0)\,T_0'(x_0)
\]
\section*{2.4  流体微团运动分析}\addcontentsline{toc}{section}{2.4 流体微团运动分析}

连续介质受力后要发生变形，流体运动过程中，流场各个微团除了移动和转动外，将发生持续的变形，而流体微团的变形状态和它的应力状态密切联系，为此我们需要对流体微团变形过程做仔细的研究.

\subsection*{1. 柯西-亥姆霍兹(Cauchy-Helmholtz)流体微团速度分解定理}

\noindent\textbf{定理2.1}\quad 流场$\boldsymbol{U}(\boldsymbol{x},t)$中微团上任意一点的运动可以分解为平动、转动和变形部分之和.

\noindent\textbf{证明}\quad 设有流场$\boldsymbol{U}(\boldsymbol{x},t)$，在时刻$t$任取一微团，微团上参考点$\boldsymbol{x}_0$处的速度为$\boldsymbol{U}_0=\boldsymbol{U}(\boldsymbol{x}_0,t)$，我们来考察微团上同一时刻任意一点$\boldsymbol{x}=\boldsymbol{x}_0+\delta\boldsymbol{x}$的速度$\boldsymbol{U}(\boldsymbol{x}_0+\delta\boldsymbol{x},t)$.利用泰勒展开，有
\[
\boldsymbol{U}(\boldsymbol{x}_0+\delta\boldsymbol{x},t)=\boldsymbol{U}(\boldsymbol{x}_0,t)+\frac{\partial\boldsymbol{U}}{\partial x_1}\delta x_1+\frac{\partial\boldsymbol{U}}{\partial x_2}\delta x_2+\frac{\partial\boldsymbol{U}}{\partial x_3}\delta x_3+O(|\delta\boldsymbol{x}|^{2})
\]
略去高阶小量后，微团上任意点速度为
\[
\boldsymbol{U}(\boldsymbol{x},t)=\boldsymbol{U}_0+\Big(\frac{\partial\boldsymbol{U}}{\partial x_j}\Big)_{\!0}\delta x_j
\]
进一步写成速度分量形式为
\[
U_i(\boldsymbol{x},t)=U_{0i}+\Big(\frac{\partial U_i}{\partial x_j}\Big)_{\!0}\delta x_j
\]
利用张量识别定理(见附录)可以证明$\dfrac{\partial U_i}{\partial x_j}$是二阶张量(证明略)，它可以进一步分解为一个对称张量和另一个反对称张量之和：
\begin{equation*}
\frac{\partial U_i}{\partial x_j}=\frac{1}{2}\Big(\frac{\partial U_i}{\partial x_j}+\frac{\partial U_j}{\partial x_i}\Big)+\frac{1}{2}\Big(\frac{\partial U_i}{\partial x_j}-\frac{\partial U_j}{\partial x_i}\Big)
\tag{2.19}
\end{equation*}
式(2.19)右端第一项用$S_{ij}$表示，它是对称张量，因为
\begin{equation*}
S_{ij}=\frac{1}{2}\Big(\frac{\partial U_i}{\partial x_j}+\frac{\partial U_j}{\partial x_i}\Big)=\frac{1}{2}\Big(\frac{\partial U_j}{\partial x_i}+\frac{\partial U_i}{\partial x_j}\Big)=S_{ji}
\tag{2.20a}
\end{equation*}
由于对称性，$S_{ij}$只有六个独立分量.

式(2.20)右端第二项用$\Omega_{ij}$表示，它是反对称张量，因为
\begin{equation*}
\Omega_{ij}=\frac{1}{2}\Big(\frac{\partial U_i}{\partial x_j}-\frac{\partial U_j}{\partial x_i}\Big)=-\frac{1}{2}\Big(\frac{\partial U_j}{\partial x_i}-\frac{\partial U_i}{\partial x_j}\Big)=-\Omega_{ji}
\tag{2.20b}
\end{equation*}
由于反对称性，$\Omega_{ij}=-\Omega_{ji}$，因此$\Omega_{\alpha\alpha}=0$(不求和)，就是说$\Omega_{ij}$只有三个独立分量.

在直角坐标系中，$S_{ij}$、$\Omega_{ij}$的各个分量表示式如下：
\begin{gather*}
S_{xx}=\frac{\partial u}{\partial x},\qquad S_{yy}=\frac{\partial v}{\partial y},\qquad S_{zz}=\frac{\partial w}{\partial z}\tag{2.21a}\\[4pt]
\begin{cases}
S_{xy}=S_{yx}=\dfrac{1}{2}\Big(\dfrac{\partial u}{\partial y}+\dfrac{\partial v}{\partial x}\Big)\\[6pt]
S_{yz}=S_{zy}=\dfrac{1}{2}\Big(\dfrac{\partial v}{\partial z}+\dfrac{\partial w}{\partial y}\Big)\\[6pt]
S_{zx}=S_{xz}=\dfrac{1}{2}\Big(\dfrac{\partial w}{\partial x}+\dfrac{\partial u}{\partial z}\Big)
\end{cases}\tag{2.21b}\\[10pt]
\Omega_{xx}=\Omega_{yy}=\Omega_{zz}=0\tag{2.22a}\\[4pt]
\begin{cases}
\Omega_{xy}=-\Omega_{yx}=\dfrac{1}{2}\Big(\dfrac{\partial u}{\partial y}-\dfrac{\partial v}{\partial x}\Big)\\[6pt]
\Omega_{yz}=-\Omega_{zy}=\dfrac{1}{2}\Big(\dfrac{\partial v}{\partial z}-\dfrac{\partial w}{\partial y}\Big)\\[6pt]
\Omega_{zx}=-\Omega_{xz}=\dfrac{1}{2}\Big(\dfrac{\partial w}{\partial x}-\dfrac{\partial u}{\partial z}\Big)
\end{cases}\tag{2.22b}
\end{gather*}
在后面将证明$S_{ij}$的6个分量分别是流体微团的线变形率和角变形率；$\Omega_{ij}$的3个分量是流体微团的准刚体转动角速度的3个分量.利用上述两个张量，流体微团上任意一点的速度可表示为
\begin{equation*}
U_i=U_{i0}+(S_{ij})_0\,\delta x_j+(\Omega_{ij})_0\,\delta x_j
\tag{2.23}
\end{equation*}
下面先证明$S_{ij}$的各个分量是流体微团的变形率，然后证明$\Omega_{ij}$是微团的准刚体角速度.完成以上证明后，定理2.1证毕.

\subsection*{2. 流体微团的变形率张量}

为了使流体微团变形的分析直观和明了，在直角坐标系中来研究流体微团的运动.在时刻$t$，取一长方体流体微团，参考点在坐标原点，长方体$OABC$在$\delta t$时刻后运动到$O'A'B'C'$，它变形为近似平行六面体(参见图2.8).也就是说，微元长方体各边$OA$，$OB$，$OC$可能伸长或缩短，微元体的矩形表面$OA\times OB$等可能因剪切而变为平行四边形$O'A'\times O'B'$等.下面将导出微元体变形的解析表达式.

先导出微元体上各点$(O,A,B,C)$的位移$\overrightarrow{OO'}$，$\overrightarrow{AA'}$，$\overrightarrow{BB'}$，$\overrightarrow{CC'}$，以及$\overrightarrow{O'A'}$，$\overrightarrow{O'B'}$，$\overrightarrow{O'C'}$的微元向量表达式(参考图2.8)，已知
\begin{equation*}
\overrightarrow{OA}=\delta x\,\boldsymbol{e}_x,\qquad \overrightarrow{OB}=\delta y\,\boldsymbol{e}_y,\qquad \overrightarrow{OC}=\delta z\,\boldsymbol{e}_z
\tag{2.24}
\end{equation*}
$O,A,B,C$各位移为
\begin{equation*}
\overrightarrow{OO'}=\boldsymbol{U}_0\,\delta t,\qquad \overrightarrow{AA'}=\boldsymbol{U}_A\,\delta t,\qquad \overrightarrow{BB'}=\boldsymbol{U}_B\,\delta t,\qquad \overrightarrow{CC'}=\boldsymbol{U}_C\,\delta t
\tag{2.25}
\end{equation*}
$A$点速度可用速度场公式求出(参见图2.8)：
\[
\boldsymbol{U}_A=\boldsymbol{U}(\boldsymbol{x}_0+\delta x\,\boldsymbol{e}_x,t)=\boldsymbol{U}_0+\Big(\frac{\partial\boldsymbol{U}}{\partial x}\Big)_{\!0}\delta x+O(\delta x^{2})
\]
同理$B$，$C$点的速度为
\[
\boldsymbol{U}_B=\boldsymbol{U}_0+\Big(\frac{\partial\boldsymbol{U}}{\partial y}\Big)_{\!0}\delta y+O(\delta y^{2}),\qquad \boldsymbol{U}_C=\boldsymbol{U}_0+\Big(\frac{\partial\boldsymbol{U}}{\partial z}\Big)_{\!0}\delta z+O(\delta z^{2})
\]
由向量几何关系式(见图2.8)可求出向量$\overrightarrow{O'A'}$，$\overrightarrow{O'B'}$，$\overrightarrow{O'C'}$的表达式(均略去高阶小量，下同)，例如：
\begin{align*}
\overrightarrow{O'A'}&=\overrightarrow{OA}+\overrightarrow{AA'}-\overrightarrow{OO'}=\delta x\,\boldsymbol{e}_x+\Big[\boldsymbol{U}_0+\Big(\frac{\partial\boldsymbol{U}}{\partial x}\Big)_{\!0}\delta x\Big]\delta t-\boldsymbol{U}_0\,\delta t\\
&=\Big[\Big(\frac{\partial\boldsymbol{U}}{\partial x}\Big)_{\!0}\delta t+\boldsymbol{e}_x\Big]\delta x
\end{align*}
将$\dfrac{\partial\boldsymbol{U}}{\partial x}=\dfrac{\partial u}{\partial x}\boldsymbol{e}_x+\dfrac{\partial v}{\partial x}\boldsymbol{e}_y+\dfrac{\partial w}{\partial x}\boldsymbol{e}_z$代入上式(为了简明起见取消下标0)，得
\[
\overrightarrow{O'A'}=\Big(1+\frac{\partial u}{\partial x}\delta t\Big)\delta x\,\boldsymbol{e}_x+\frac{\partial v}{\partial x}\delta t\,\delta x\,\boldsymbol{e}_y+\frac{\partial w}{\partial x}\delta t\,\delta x\,\boldsymbol{e}_z
\]
同理可得
\begin{align*}
\overrightarrow{O'B'}&=\frac{\partial u}{\partial y}\delta t\,\delta y\,\boldsymbol{e}_x+\Big(1+\frac{\partial v}{\partial y}\delta t\Big)\delta y\,\boldsymbol{e}_y+\frac{\partial w}{\partial y}\delta t\,\delta y\,\boldsymbol{e}_z\\
\overrightarrow{O'C'}&=\frac{\partial u}{\partial z}\delta t\,\delta z\,\boldsymbol{e}_x+\frac{\partial v}{\partial z}\delta t\,\delta z\,\boldsymbol{e}_y+\Big(1+\frac{\partial w}{\partial z}\delta t\Big)\delta z\,\boldsymbol{e}_z
\end{align*}
由以上位移公式可导出变形率公式.

(1)线变形率

\noindent\textbf{定义2.7}\quad 单位时间内微元线段$\delta x_i$的相对伸长，称为线变形率，并用$\dot{\epsilon}_i$表示，下标$i$表示线段的方向.

$x$方向微元线段的伸长率应为
\begin{equation*}
\dot{\epsilon}_x=\lim_{\delta t\to 0}\frac{|\overrightarrow{O'A'}|-|\overrightarrow{OA}|}{|\overrightarrow{OA}|\,\delta t}
\tag{2.26}
\end{equation*}
已知$|\overrightarrow{OA}|=\delta x$，由上面已导出的$\overrightarrow{O'A'}$公式，不难得到
\begin{align*}
|\overrightarrow{O'A'}|&=\sqrt{\Big(1+\frac{\partial u}{\partial x}\delta t\Big)^{2}\delta x^{2}+\Big(\frac{\partial v}{\partial x}\Big)^{2}\delta t^{2}\delta x^{2}+\Big(\frac{\partial w}{\partial x}\Big)^{2}\delta t^{2}\delta x^{2}}\\
&=\delta x\sqrt{1+2\frac{\partial u}{\partial x}\delta t+\Big[\Big(\frac{\partial u}{\partial x}\Big)^{2}+\Big(\frac{\partial v}{\partial x}\Big)^{2}+\Big(\frac{\partial w}{\partial x}\Big)^{2}\Big]\delta t^{2}}
\end{align*}
利用二项式展开，上式等于
\[
|\overrightarrow{O'A'}|=\delta x\Big[1+\frac{\partial u}{\partial x}\delta t+O(\delta t)^{2}\Big]
\]
将它代入式(2.26)，得
\begin{gather*}
\dot{\epsilon}_x=\frac{\partial u}{\partial x}\tag{2.27a}\\[2pt]
\dot{\epsilon}_y=\frac{\partial v}{\partial y}\tag{2.27b}\\[2pt]
\dot{\epsilon}_z=\frac{\partial w}{\partial z}\tag{2.27c}
\end{gather*}
于是证明了直角坐标系中$S_{ij}$张量的对角线分量分别表示所在方向的线变形率.

(2)角变形率

\noindent\textbf{定义2.8}\quad 微元平面上两垂直线段夹角在单位时间内减小量之半称为该面的角变形率，并用$\dot{\gamma}_{ij}$表示，下标表示两线段所在平面.例如$\dot{\gamma}_{xy}$表示$x\text{-}y$平面上的角变形率.

按定义$x\text{-}y$平面上切变率$\dot{\gamma}_{xy}$公式应是：
\begin{equation*}
\dot{\gamma}_{xy}=\lim_{\delta t\to 0}\frac{\angle AOB-\angle A'O'B'}{2\delta t}
\tag{2.28}
\end{equation*}
变形前的长方体流体微团中$\angle AOB=\dfrac{\pi}{2}$，在$\delta t$时刻内$\angle AOB-\angle A'O'B'$是小量，可以用它的正弦来量度，
\begin{align*}
\angle AOB-\angle A'O'B'&\approx\sin(\angle AOB-\angle A'O'B')\\
&=\sin\Big(\frac{\pi}{2}-\angle A'O'B'\Big)=\cos(\angle A'O'B')
\end{align*}
应用两向量点积公式：$\boldsymbol{X}\cdot\boldsymbol{Y}=|\boldsymbol{X}||\boldsymbol{Y}|\cos(\boldsymbol{X},\boldsymbol{Y})$，可得
\[
\cos(\angle A'O'B')=\frac{\overrightarrow{O'A'}\cdot\overrightarrow{O'B'}}{|\overrightarrow{O'A'}|\,|\overrightarrow{O'B'}|}
\]
前面已经导出
\begin{align*}
\overrightarrow{O'A'}&=\Big(1+\frac{\partial u}{\partial x}\delta t\Big)\delta x\,\boldsymbol{e}_x+\frac{\partial v}{\partial x}\delta x\,\delta t\,\boldsymbol{e}_y+\frac{\partial w}{\partial x}\delta x\,\delta t\,\boldsymbol{e}_z\\
\overrightarrow{O'B'}&=\frac{\partial u}{\partial y}\delta y\,\delta t\,\boldsymbol{e}_x+\Big(1+\frac{\partial v}{\partial y}\delta t\Big)\delta y\,\boldsymbol{e}_y+\frac{\partial w}{\partial y}\delta y\,\delta t\,\boldsymbol{e}_z
\end{align*}
以上两向量的点积等于：
\[
\overrightarrow{O'A'}\cdot\overrightarrow{O'B'}=\frac{\partial u}{\partial y}\delta t\,\delta x\,\delta y+\frac{\partial v}{\partial x}\delta t\,\delta x\,\delta y+O(\delta t^{2}\delta x\,\delta y)
\]
$\overrightarrow{O'A'}$，$\overrightarrow{O'B'}$的模分别为
\begin{gather*}
|\overrightarrow{O'A'}|=\Big(1+\frac{\partial u}{\partial x}\delta t\Big)\delta x+O(\delta t^{2})\\[2pt]
|\overrightarrow{O'B'}|=\Big(1+\frac{\partial v}{\partial y}\delta t\Big)\delta y+O(\delta t^{2})
\end{gather*}
代入$\cos(\angle A'O'B')$得
\begin{align*}
\sin(\angle AOB-\angle A'O'B')=\cos(\angle A'O'B')&=\frac{\Big(\dfrac{\partial u}{\partial y}+\dfrac{\partial v}{\partial x}\Big)\delta t\,\delta x\,\delta y}{\delta x\,\delta y\,[1+O(\delta t)]}\\
&=\Big(\frac{\partial u}{\partial y}+\frac{\partial v}{\partial x}\Big)\delta t+O(\delta t^{2})
\end{align*}
于是
\[
\angle AOB-\angle A'O'B'=\Big(\frac{\partial u}{\partial y}+\frac{\partial v}{\partial x}\Big)\delta t+O(\delta t^{2})
\]
代入角变形率定义式，得
\begin{gather*}
\dot{\gamma}_{xy}=\frac{1}{2}\Big(\frac{\partial u}{\partial y}+\frac{\partial v}{\partial x}\Big)\tag{2.29a}\\[2pt]
\dot{\gamma}_{yz}=\frac{1}{2}\Big(\frac{\partial v}{\partial z}+\frac{\partial w}{\partial y}\Big)\tag{2.29b}\\[2pt]
\dot{\gamma}_{zx}=\frac{1}{2}\Big(\frac{\partial w}{\partial x}+\frac{\partial u}{\partial z}\Big)\tag{2.29c}
\end{gather*}
于是证明了，二阶对称张量的另外3个分量分别表示角变形率.综合以上结果，对称张量$S_{ij}$的6个分量分别是微团的线变形率和角变形率：
\begin{equation*}
S_{ij}=\begin{bmatrix}
\dot{\epsilon}_1 & \dot{\gamma}_{12} & \dot{\gamma}_{13}\\
\dot{\gamma}_{12} & \dot{\epsilon}_2 & \dot{\gamma}_{23}\\
\dot{\gamma}_{13} & \dot{\gamma}_{23} & \dot{\epsilon}_3
\end{bmatrix}
\tag{2.30}
\end{equation*}
因此速度梯度张量的对称部分$S_{ij}$又称变形率张量.

(3)体积变化率

\noindent\textbf{定义2.9}\quad 流体微团的体积变化率是单位时间内流体微团体积的相对增长率，并用$\dot{\epsilon}$表示.

按定义
\begin{equation*}
\dot{\epsilon}=\lim_{\delta t\to 0}\frac{V(O'A'B'C')-V(OABC)}{V(OABC)\,\delta t}
\tag{2.31}
\end{equation*}
可以证明体积变化率有以下公式：
\[
\dot{\epsilon}=\dot{\epsilon}_x+\dot{\epsilon}_y+\dot{\epsilon}_z=S_{11}+S_{22}+S_{33}
\]
或写成
\begin{equation*}
\dot{\epsilon}=S_{ii}
\tag{2.32}
\end{equation*}
式(2.32)证明如下：由3个向量$\boldsymbol{X},\boldsymbol{Y},\boldsymbol{Z}$组成的平行六面体的体积为
\[
V=\boldsymbol{X}\cdot(\boldsymbol{Y}\times\boldsymbol{Z})=\begin{vmatrix}
X_1 & X_2 & X_3\\
Y_1 & Y_2 & Y_3\\
Z_1 & Z_2 & Z_3
\end{vmatrix}
\]
于是体积$V(OABC)=\overrightarrow{OA}\cdot(\overrightarrow{OB}\times\overrightarrow{OC})=\delta x\,\delta y\,\delta z$.体积$V(O'A'B'C')=\overrightarrow{O'A'}\cdot(\overrightarrow{O'B'}\times\overrightarrow{O'C'})$，将$\overrightarrow{O'A'}$，$\overrightarrow{O'B'}$，$\overrightarrow{O'C'}$代入，得
\begin{align*}
\overrightarrow{O'A'}\cdot(\overrightarrow{O'B'}\times\overrightarrow{O'C'})&=\begin{vmatrix}
\Big(1+\dfrac{\partial u}{\partial x}\delta t\Big) & \dfrac{\partial v}{\partial x}\delta t & \dfrac{\partial w}{\partial x}\delta t\\[8pt]
\dfrac{\partial u}{\partial y}\delta t & \Big(1+\dfrac{\partial v}{\partial y}\delta t\Big) & \dfrac{\partial w}{\partial y}\delta t\\[8pt]
\dfrac{\partial u}{\partial z}\delta t & \dfrac{\partial v}{\partial z}\delta t & \Big(1+\dfrac{\partial w}{\partial z}\delta t\Big)
\end{vmatrix}\delta x\,\delta y\,\delta z\\[8pt]
&=\Big[\Big(1+\frac{\partial u}{\partial x}\delta t\Big)\Big(1+\frac{\partial v}{\partial y}\delta t\Big)\Big(1+\frac{\partial w}{\partial z}\delta t\Big)+O(\delta t^{2})\Big]\delta x\,\delta y\,\delta z\\
&=\Big[1+\Big(\frac{\partial u}{\partial x}+\frac{\partial v}{\partial y}+\frac{\partial w}{\partial z}\Big)\delta t\Big]\delta x\,\delta y\,\delta z+O(\delta t^{2})\,\delta x\,\delta y\,\delta z
\end{align*}
代入体积变化率定义得式(2.31)
\[
\dot{\epsilon}=\frac{\partial u}{\partial x}+\frac{\partial v}{\partial y}+\frac{\partial w}{\partial z}=S_{ii}
\]
应用向量场散度公式$\nabla\cdot\boldsymbol{a}=\dfrac{\partial a_i}{\partial x_i}$(见附录)，直角坐标系中速度场的散度等于：
\[
\nabla\cdot\boldsymbol{U}=\frac{\partial u}{\partial x}+\frac{\partial v}{\partial y}+\frac{\partial w}{\partial z}
\]
因此体积变化率公式(2.32)又可写作：
\begin{equation*}
\dot{\epsilon}=\nabla\cdot\boldsymbol{U}
\tag{2.33}
\end{equation*}
式(2.32)和式(2.33)表明，微元体的体积增长率等于变形率张量对角线上3项之和(又称张量的迹)，也等于该点速度场散度，它是一个标量场.

当流体微团为不可压缩时，体积增长率等于零，也就是该点的速度场散度等于零.

从变形率张量中减去各向同性的体膨胀率称之为纯变形率张量：
\[
S_{ij}'=S_{ij}-\frac{1}{3}\dot{\epsilon}\,\delta_{ij}
\]
在正交曲线坐标系中，沿曲线坐标的线变形率和坐标曲面上的角变形率的计算公式可用向量导数方法求出.具体做法如下：首先将直角坐标系导出的变形率公式用一般向量运算表达为
\begin{equation*}
S_{ij}=\frac{1}{2}\Big(\frac{\partial u_i}{\partial x_j}+\frac{\partial u_j}{\partial x_i}\Big)=\frac{1}{2}\big[(\nabla\boldsymbol{U})+(\nabla\boldsymbol{U})^{\mathrm{T}}\big]_{ij}
\tag{2.34}
\end{equation*}
式中$\nabla\boldsymbol{U}=\dfrac{\partial u_i}{\partial x_j}\boldsymbol{e}_i\boldsymbol{e}_j$，$(\nabla\boldsymbol{U})^{\mathrm{T}}=\dfrac{\partial u_j}{\partial x_i}\boldsymbol{e}_i\boldsymbol{e}_j$是速度梯度张量$\nabla\boldsymbol{U}$的转置张量.因为张量表达式和坐标系无关，所以$S_{ij}=\big[(\nabla\boldsymbol{U})+(\nabla\boldsymbol{U})^{\mathrm{T}}\big]_{ij}/2$在任意正交坐标系中都成立.第二步，将算符$\nabla$和速度场$\boldsymbol{U}$都用曲线坐标表示，然后按向量求导法则分别求出$(\nabla\boldsymbol{U})^{\mathrm{T}}$和$(\nabla\boldsymbol{U})$，代入式(2.34)就可得到曲线坐标系中的变形率公式.下面举例说明运算过程.

\noindent\textbf{例2.7}\quad 给定平面流场的极坐标表达式：$v_r=u(r,\theta)$，$v_\theta=v(r,\theta)$，求流动平面上向径方向和圆周方向的线变形率，以及平面上的角变形率.

\noindent\textbf{解}\quad 将速度和向量算符写成极坐标式：$\boldsymbol{U}=u\,\boldsymbol{e}_r+v\,\boldsymbol{e}_\theta$，$\nabla=\boldsymbol{e}_r\dfrac{\partial}{\partial r}+\boldsymbol{e}_\theta\dfrac{\partial}{r\partial\theta}$，然后计算$\nabla\boldsymbol{U}$：
\begin{align*}
\nabla\boldsymbol{U}&=\Big(\boldsymbol{e}_r\frac{\partial}{\partial r}+\boldsymbol{e}_\theta\frac{\partial}{r\partial\theta}\Big)(u\,\boldsymbol{e}_r+v\,\boldsymbol{e}_\theta)
=\boldsymbol{e}_r\frac{\partial}{\partial r}(u\,\boldsymbol{e}_r+v\,\boldsymbol{e}_\theta)+\boldsymbol{e}_\theta\frac{\partial}{r\partial\theta}(u\,\boldsymbol{e}_r+v\,\boldsymbol{e}_\theta)\\
&=\boldsymbol{e}_r\frac{\partial u}{\partial r}\boldsymbol{e}_r+\boldsymbol{e}_r u\frac{\partial\boldsymbol{e}_r}{\partial r}+\boldsymbol{e}_r\frac{\partial v}{\partial r}\boldsymbol{e}_\theta+\boldsymbol{e}_r v\frac{\partial\boldsymbol{e}_\theta}{\partial r}
+\boldsymbol{e}_\theta\frac{\partial u}{r\partial\theta}\boldsymbol{e}_r+\boldsymbol{e}_\theta u\frac{\partial\boldsymbol{e}_r}{r\partial\theta}
+\boldsymbol{e}_\theta\frac{\partial v}{r\partial\theta}\boldsymbol{e}_\theta+\boldsymbol{e}_\theta v\frac{\partial\boldsymbol{e}_\theta}{r\partial\theta}
\end{align*}
径向和周向单位向量的导数分别有公式(参见附录)：
\[
\frac{\partial\boldsymbol{e}_r}{\partial r}=0,\qquad \frac{\partial\boldsymbol{e}_\theta}{\partial r}=0,\qquad \frac{\partial\boldsymbol{e}_r}{\partial\theta}=\boldsymbol{e}_\theta,\qquad \frac{\partial\boldsymbol{e}_\theta}{\partial\theta}=-\boldsymbol{e}_r
\]
于是有
\[
\nabla\boldsymbol{U}=\begin{pmatrix}
\dfrac{\partial u}{\partial r}\boldsymbol{e}_r\boldsymbol{e}_r & \dfrac{\partial v}{\partial r}\boldsymbol{e}_r\boldsymbol{e}_\theta\\[10pt]
\Big(\dfrac{\partial u}{r\partial\theta}-\dfrac{v}{r}\Big)\boldsymbol{e}_\theta\boldsymbol{e}_r & \Big(\dfrac{\partial v}{r\partial\theta}+\dfrac{u}{r}\Big)\boldsymbol{e}_\theta\boldsymbol{e}_\theta
\end{pmatrix}
\ \text{以及}\
(\nabla\boldsymbol{U})^{\mathrm{T}}=\begin{pmatrix}
\dfrac{\partial u}{\partial r}\boldsymbol{e}_r\boldsymbol{e}_r & \Big(\dfrac{\partial u}{r\partial\theta}-\dfrac{v}{r}\Big)\boldsymbol{e}_\theta\boldsymbol{e}_r\\[10pt]
\dfrac{\partial v}{\partial r}\boldsymbol{e}_r\boldsymbol{e}_\theta & \Big(\dfrac{\partial v}{r\partial\theta}+\dfrac{u}{r}\Big)\boldsymbol{e}_\theta\boldsymbol{e}_\theta
\end{pmatrix}
\]
因而线变形率等于：
\[
\dot{\epsilon}_r=\frac{1}{2}(\nabla\boldsymbol{U}+(\nabla\boldsymbol{U})^{\mathrm{T}})_{rr}=\frac{\partial u}{\partial r},\qquad
\dot{\epsilon}_\theta=\frac{1}{2}(\nabla\boldsymbol{U}+(\nabla\boldsymbol{U})^{\mathrm{T}})_{\theta\theta}=\frac{\partial v}{r\partial\theta}+\frac{u}{r}
\]
角变形率等于：
\[
\dot{\gamma}_{r\theta}=\dot{\gamma}_{\theta r}=\frac{1}{2}(\nabla\boldsymbol{U}+(\nabla\boldsymbol{U})^{\mathrm{T}})_{\theta r}=\frac{1}{2}\Big(\frac{\partial u}{r\partial\theta}+\frac{\partial v}{\partial r}-\frac{v}{r}\Big)
\]
附录中给出了常用正交曲线坐标系中的变形率公式，读者可以选择几个公式作为向量运算的练习.

\noindent\textbf{例2.8}\quad 给定直角坐标系中速度场：$u=x^{2}y+y^{2}$，$v=x^{2}-xy^{2}$，$w=0$，求各变形率张量的分量，并判断该流场是否为不可压缩流场.

\noindent\textbf{解}\quad 变形率张量为
\begin{gather*}
\dot{\epsilon}_x=\frac{\partial u}{\partial x}=2xy,\qquad \dot{\epsilon}_y=\frac{\partial v}{\partial y}=-2xy,\qquad \dot{\epsilon}_z=\frac{\partial w}{\partial z}=0\\[4pt]
\dot{\gamma}_{xy}=\frac{1}{2}\Big(\frac{\partial u}{\partial y}+\frac{\partial v}{\partial x}\Big)=\frac{x^{2}-y^{2}}{2}+x+y\\[4pt]
\dot{\gamma}_{yz}=\frac{1}{2}\Big(\frac{\partial v}{\partial z}+\frac{\partial w}{\partial y}\Big)=0,\qquad
\dot{\gamma}_{zx}=\frac{1}{2}\Big(\frac{\partial w}{\partial x}+\frac{\partial u}{\partial z}\Big)=0
\end{gather*}
计算散度$\nabla\cdot\boldsymbol{U}$可判断是否为不可压缩流场，因有
\[
\nabla\cdot\boldsymbol{U}=\frac{\partial u}{\partial x}+\frac{\partial v}{\partial y}+\frac{\partial w}{\partial z}=\dot{\epsilon}_x+\dot{\epsilon}_y+\dot{\epsilon}_z=2xy-2xy=0
\]
故该流场为不可压缩流场.
\subsection*{3. 流体微团的准刚体转动角速度}

现在研究二阶反对称张量$\Omega_{ij}$.可以证明反对称张量可以和一个向量对等.例如，$\Omega_{ij}$的直角坐标表达式为
\begin{equation*}
\Omega_{ij}=\begin{bmatrix}
0 & \dfrac{1}{2}\Big(\dfrac{\partial u}{\partial y}-\dfrac{\partial v}{\partial x}\Big) & \dfrac{1}{2}\Big(\dfrac{\partial u}{\partial z}-\dfrac{\partial w}{\partial x}\Big)\\[12pt]
\dfrac{1}{2}\Big(\dfrac{\partial v}{\partial x}-\dfrac{\partial u}{\partial y}\Big) & 0 & \dfrac{1}{2}\Big(\dfrac{\partial v}{\partial z}-\dfrac{\partial w}{\partial y}\Big)\\[12pt]
\dfrac{1}{2}\Big(\dfrac{\partial w}{\partial x}-\dfrac{\partial u}{\partial z}\Big) & \dfrac{1}{2}\Big(\dfrac{\partial w}{\partial y}-\dfrac{\partial v}{\partial z}\Big) & 0
\end{bmatrix}
\tag{2.35}
\end{equation*}
另外，速度场旋度$(\boldsymbol{\omega}=\nabla\times\boldsymbol{U})$在直角坐标系中的表达式为
\begin{equation*}
\omega_x=\Big(\frac{\partial w}{\partial y}-\frac{\partial v}{\partial z}\Big),\qquad \omega_y=\Big(\frac{\partial u}{\partial z}-\frac{\partial w}{\partial x}\Big),\qquad \omega_z=\Big(\frac{\partial v}{\partial x}-\frac{\partial u}{\partial y}\Big)
\tag{2.36}
\end{equation*}
和式(2.35)对比，可以将反对称张量表示为
\[
2\Omega_{ij}=\begin{bmatrix}
0 & -\omega_z & \omega_y\\
\omega_z & 0 & -\omega_x\\
-\omega_y & \omega_x & 0
\end{bmatrix}
\]
或可写成
\begin{equation*}
2\Omega_{ij}=-\epsilon_{ijk}\,\omega_k
\tag{2.37}
\end{equation*}
式中$\epsilon_{ijk}$是置换符号，这就证明了反对称张量$\Omega_{ij}$和速度场旋度之半是等价的.进一步考察反对称张量$\Omega_{ij}$在流体微团速度分解中的贡献，将式(2.37)代入式(2.23)中的$(\Omega_{ij})_0\delta x_j$，得
\begin{equation*}
\Omega_{ij}\,\delta x_j=-\frac{\epsilon_{ijk}\,\omega_k\,\delta x_j}{2}
\tag{2.38}
\end{equation*}
利用向量叉积公式(见附录)：$(\boldsymbol{a}\times\boldsymbol{b})_i=\epsilon_{ijk}a_jb_k$，式(2.38)可写作：
\[
\Omega_{ij}\,\delta x_j=-\frac{\epsilon_{ijk}\,\omega_k\,\delta x_j}{2}=\frac{(-\delta\boldsymbol{x}\times\boldsymbol{\omega})_i}{2}=\frac{(\boldsymbol{\omega}\times\delta\boldsymbol{x})_i}{2}
\]
在刚体运动学中我们知道刚体以角速度$\boldsymbol{\Omega}_r$旋转时，刚体上任意点相对于参考点的线速度为$(\boldsymbol{\Omega}_r\times\delta\boldsymbol{r})_i=-\epsilon_{ijk}\Omega_{rk}\,\delta x_j$.对照式(2.38)，旋度之半相当于刚体角速度.这就是说反对称张量$\Omega_{ij}$在流体微团速度分解中的贡献和微团刚体转动相当.定义$\Omega_{ij}$(也就是$\dfrac{\boldsymbol{\omega}}{2}$)为流体微团的\textbf{准刚体角速度}.由于不均匀流场中流体不断变形，一般来说，在不均匀流场中微团的准刚体角速度的大小和方向随空间位置而改变.不仅如此，在非定常流场中，准刚体角速度还随时间变化.

流体微团准刚体转动应理解为微团在当时当地绕微团中心的刚体转动，形象地说，流体微团的准刚体转动就像地球的自转一样，它是绕自身某一轴线转动.流体微团作为一个实体还可绕某点作圆周运动，好比地球绕太阳的公转.必须把微团的准刚体转动和质点的圆周运动区别开来.下面举例来说明这一概念.

\noindent\textbf{例2.9}\quad 给定两个流场：$(1)\,u=-y,\ v=x,\ w=0$；$(2)\,u=-\dfrac{y}{x^{2}+y^{2}},\ v=\dfrac{x}{x^{2}+y^{2}},\ w=0$.求这两个流场的迹线和准刚体角速度.

\noindent\textbf{解}\quad 两流场都是平面定常流动，只须求$x\text{-}y$平面上的流线，它们同时也是迹线.应用流线微分方程
\[
\frac{\mathrm{d}x}{u}=\frac{\mathrm{d}y}{v}
\]
将两速度场公式分别代入，得到相同的流线方程：
\[
\frac{\mathrm{d}x}{-y}=\frac{\mathrm{d}y}{x}
\]
积分后得
\[
x^{2}+y^{2}=c
\]
上式表明：两流场的流体微团迹线都是同心圆，即任一流体微团都作圆周运动.

利用准刚体角速度公式：
\[
\omega_x=\frac{1}{2}\Big(\frac{\partial w}{\partial y}-\frac{\partial v}{\partial z}\Big),\qquad \omega_y=\frac{1}{2}\Big(\frac{\partial u}{\partial z}-\frac{\partial w}{\partial x}\Big),\qquad \omega_z=\frac{1}{2}\Big(\frac{\partial v}{\partial x}-\frac{\partial u}{\partial y}\Big)
\]
对于流场(1)\quad $u=-y$，$v=x$，$w=0$，$\omega_z=1$，$\omega_x=0$，$\omega_y=0$；

对于流场(2)\quad $u=\dfrac{-y}{x^{2}+y^{2}}$，$v=\dfrac{x}{x^{2}+y^{2}}$，$w=0$，$\omega_z=0$，$\omega_x=0$，$\omega_y=0$.

以上分析表明：虽然两个流场的质点都作圆周运动，但是流场(1)中流体微团的准刚体转动角速度处处为1，就是说流体微团在作圆周运动的同时还有自转；对于流场(2)，流体微团的准刚体转动角速度处处为零，也就是说流体微团作圆周运动而没有自转.

为了形象地解释这两个流场，可以把流体微团做上标记，以考察流体微团在运动过程中的“自转”，即准刚体转动，如图2.9所示.对于流场(1)，流体微团作圆周运动的同时以刚体转动的方式绕微团中心转动(矩形的黑色标记随着微团转动)；流场(2)流体微团只有圆周平移运动，而自身没有转动(黑色标记没有转动).

为了更好地理解上面用分析方法导出的流体微团的角变形和旋转，在平面上考察一个微团的运动(图2.10).

在前面分析中已经导出，$A'$相对于$O'$的位移等于$\Big(\dfrac{\partial v}{\partial x}\Big)\delta x\,\mathrm{d}t$，因此$OA$边到$O'A'$边的转角等于$\beta=\sin\beta=\Big(\dfrac{\partial v}{\partial x}\Big)\mathrm{d}t$，同理，$OB$到$O'B'$的转角$\alpha=\sin\alpha=\Big(\dfrac{\partial u}{\partial y}\Big)\mathrm{d}t$(注意$\mathrm{d}t$是小量).于是由$\angle AOB=\dfrac{\pi}{2}$变形到$\angle A'O'B'$时，其夹角减少了$\alpha+\beta=\Big(\dfrac{\partial u}{\partial y}+\dfrac{\partial v}{\partial x}\Big)\mathrm{d}t$，它的一半就是角变形率$S_{xy}=\dfrac{\dfrac{\partial u}{\partial y}+\dfrac{\partial v}{\partial x}}{2}$和$\mathrm{d}t$的乘积.再考察平面微团的对角线$AD$的转角$\gamma$(以逆时针为正)，它应等于
\[
\gamma=\frac{1}{2}(\beta-\alpha)=\frac{1}{2}\Big(\frac{\partial v}{\partial x}-\frac{\partial u}{\partial y}\Big)\mathrm{d}t=\frac{\omega_z}{2}\mathrm{d}t
\]
也就是说，微团对角线的角速度等于微团准刚体角速度(旋度之半).微团的另外两个坐标面上的角变形率和角速度可以用同样的几何分析方法来考察.

最后回到式(2.24)，已经证明了$S_{ij}$为流体微团的变形率张量，$\Omega_{ij}$为流体微团的准刚体角速度，于是完成柯西-亥姆霍兹定理的证明.综合以上分析，速度分解定理可表述如下：

\textbf{流体微团上任意点的运动$=$参考点的移动$+$体积膨胀运动$+$纯变形运动$+$准刚体转动}

\section*{2.5  流场的旋度}\addcontentsline{toc}{section}{2.5 流场的旋度}

\subsection*{1. 涡量场及其性质}

\noindent\textbf{定义2.10}\quad 速度场的旋度$\nabla\times\boldsymbol{U}$称为涡量，并用$\boldsymbol{\omega}$表示.

根据前面式(2.36)，流体微团的准刚体角速度是速度场的涡量之半，因此，速度场旋度的性质就是流体微团准刚体运动的性质.

(1)涡量场的散度等于零

由向量运算公式，很容易证明
\[
\nabla\cdot(\nabla\times\boldsymbol{U})=0
\]
即
\begin{equation*}
\nabla\cdot\boldsymbol{\omega}=0
\tag{2.39}
\end{equation*}

(2)涡线、涡面、涡管

类似于流场的几何描述，涡量场也可以用涡线、涡面和涡管来描述.

\noindent\textbf{定义2.11}\quad 涡量场的向量线称为涡线.

类似于流线方程，涡线方程为(时间$t$为参量)
\begin{equation*}
\frac{\mathrm{d}x}{\omega_x}=\frac{\mathrm{d}y}{\omega_y}=\frac{\mathrm{d}z}{\omega_z}
\tag{2.40}
\end{equation*}
在给定时刻式(2.40)的积分曲线称为涡线.涡线的切线和当地的涡量或准刚体角速度重合，所以涡线是流体微团准刚体转动方向的连线，形象地说，涡线像一根柔性轴把流体微团穿在一起.

类似于流面和流管，也常用涡面或涡管来描述涡量场.

\noindent\textbf{定义2.12}\quad 给定瞬间，通过某一曲线(本身不是涡线)的所有涡线构成的曲面称为涡面.

\noindent\textbf{定义2.13}\quad 管状涡面的内域称为涡管.

根据涡面的定义，涡面对于涡量具有不可穿透性，即在涡面上有
\begin{equation*}
\boldsymbol{\omega}\cdot\boldsymbol{n}=0
\tag{2.41}
\end{equation*}

(3)涡通量与涡管强度

给定空间曲面，则下列面积分称为涡通量
\begin{equation*}
I=\iint_A \boldsymbol{\omega}\cdot\boldsymbol{n}\,\mathrm{d}A
\tag{2.42}
\end{equation*}
$\boldsymbol{n}$是曲面$A$的外法线方向，$I>0$，涡通量是正值；$I<0$，涡通量是负值.涡管截面上涡通量的绝对值$|I|$称为涡管强度.

\noindent\textbf{定理2.2(涡管强度守恒)}\quad 给定瞬间涡管任意截面上的涡管强度相等.

\noindent\textbf{证明}\quad 在涡管上任取两个截面$A_1$，$A_2$(图2.11)，由于涡量场散度等于零，在$A_1$，$A_2$和涡管表面$A$构成的内域有
\[
\nabla\cdot\boldsymbol{\omega}=0
\]
从而有
\[
\iiint_V \nabla\cdot\boldsymbol{\omega}\,\mathrm{d}V=0
\]
利用高斯公式
\[
\iiint_V \nabla\cdot\boldsymbol{\omega}\,\mathrm{d}V=\oint_{A_1+A_2+A} \boldsymbol{n}\cdot\boldsymbol{\omega}\,\mathrm{d}A
\]
于是有
\[
\iint_{A_1}\boldsymbol{n}\cdot\boldsymbol{\omega}\,\mathrm{d}A+\iint_{A_2}\boldsymbol{n}\cdot\boldsymbol{\omega}\,\mathrm{d}A+\iint_{A}\boldsymbol{n}\cdot\boldsymbol{\omega}\,\mathrm{d}A=0
\]
由涡管定义可知涡管侧面上$\boldsymbol{n}\cdot\boldsymbol{\omega}=0$，故有
\[
\iint_{A_1}\boldsymbol{n}\cdot\boldsymbol{\omega}\,\mathrm{d}A=-\iint_{A_2}\boldsymbol{n}\cdot\boldsymbol{\omega}\,\mathrm{d}A
\]
对上式两边分别取绝对值，证明了涡管任意截面上的涡管强度相等.

由涡管强度守恒定理可得出推论：\textbf{在流场中涡管不能消失}.假定涡管截面在某处收缩到零，如图2.12(a)，则由涡管强度守恒定理可以证明：该处的涡量必为无穷大，这在实际中是不存在的.如果涡管在某截面处突然中断，即涡管存在一个端面，在该端面以外，涡量突然等于零.由涡管强度守恒定理可以证明，这种情况也是不可能的.我们可以围绕突然中断的涡管断面作一封闭曲面，则在该封闭曲面上只有涡通量进入，而无涡通量输出，这和涡量场的散度等于零相矛盾，是运动学上不允许的.涡量强度守恒定理说明：流场中涡管只能有以下三种形式(如图2.12(b)所示).

(1)涡管两端延伸到无限远；

(2)涡管形成封闭涡环；

(3)涡管中止于物面或其他界面上.

\subsection*{2. 速度环量}

\noindent\textbf{定义2.14}\quad 在速度场中沿封闭周线的线积分$\displaystyle\oint_l\boldsymbol{U}\cdot\mathrm{d}\boldsymbol{x}$称为绕该周线的速度环量，记作$\Gamma_l$.

\noindent\textbf{定理2.3}\quad 速度环量等于张在封闭周线$l$上任意曲面的涡通量，其中曲面的法向量$\boldsymbol{n}$由右手法则确定(右手四指沿周线方向，拇指指向法线方向).

\noindent\textbf{证明}\quad 由向量场的斯托克斯公式(见附录)
\[
\oint_l\boldsymbol{U}\cdot\mathrm{d}\boldsymbol{x}=\iint_A(\nabla\times\boldsymbol{U})\cdot\boldsymbol{n}\,\mathrm{d}A
\]
而$\boldsymbol{\omega}=\nabla\times\boldsymbol{U}$，因而有
\begin{equation*}
\oint_l\boldsymbol{U}\cdot\mathrm{d}\boldsymbol{x}=\iint_A \boldsymbol{\omega}\cdot\boldsymbol{n}\,\mathrm{d}A
\tag{2.43}
\end{equation*}
等式左边是速度环量，右边是张在封闭曲线$l$上曲面的涡通量，于是证毕.

如果封闭周线围绕涡管侧面，则沿封闭周线的速度环量等于涡管涡通量，因此还可以用速度环量来量度或判断涡量场.例如流场任意周线上的速度环量都等于零，则应用速度环量定理，通过任意曲面上的涡通量等于零，也就是涡量处处等于零.

\subsection*{3. 无旋流动和速度势}

\noindent\textbf{定义2.15}\quad 速度旋度处处为零的流场称为无旋流场.

就是说，如果全流场中
\[
\nabla\times\boldsymbol{U}=\boldsymbol{0}
\]
则称该速度场为无旋流动.

无旋流动的主要特性是速度场可以用一个标量函数的梯度表示，即无旋流动一定存在一势函数$\Phi(x,y,z,t)$，它的梯度等于流场速度：
\begin{equation*}
\boldsymbol{U}=\nabla\Phi
\tag{2.44}
\end{equation*}
称$\Phi$为速度势.用向量分析公式来证明这一性质，利用斯托克斯公式我们有
\[
\oint_l\boldsymbol{U}\cdot\mathrm{d}\boldsymbol{x}=\iint_A \boldsymbol{\omega}\cdot\boldsymbol{n}\,\mathrm{d}A
\]
由于流场中处处$\boldsymbol{\omega}=\boldsymbol{0}$，因此绕任意周线$l$，速度环量都等于零，即
\[
\oint_l\boldsymbol{U}\cdot\mathrm{d}\boldsymbol{x}=0
\]
因为封闭周线$l$是任意曲线，所以$\displaystyle\int\boldsymbol{U}\cdot\mathrm{d}\boldsymbol{x}$和积分路径无关.也就是说，流场中任意两点间积分$\displaystyle\int_{x_1}^{x}\boldsymbol{U}\cdot\mathrm{d}\boldsymbol{x}$只是空间点$\boldsymbol{x}$的函数，将该函数定义为势函数：
\begin{equation*}
\Phi(\boldsymbol{x},t)=\int_{x_1}^{x}\boldsymbol{U}\cdot\mathrm{d}\boldsymbol{x}
\tag{2.45}
\end{equation*}
把它写成全微分式
\[
\mathrm{d}\Phi=U_1\mathrm{d}x_1+U_2\mathrm{d}x_2+U_3\mathrm{d}x_3
\]
就可得
\[
U_1=\frac{\partial\Phi}{\partial x_1},\qquad U_2=\frac{\partial\Phi}{\partial x_2},\qquad U_3=\frac{\partial\Phi}{\partial x_3}
\]
于是证明了流体的无旋运动必存在速度势，并导出速度势和速度场之间的关系式(2.44)、式(2.45).

本书将在第4章详细讨论理想不可压缩流体的无旋流动.
\section*{2.6  给定流场的散度和旋度求速度场}\addcontentsline{toc}{section}{2.6 给定流场的散度和旋度求速度场}

从2.5节的微团运动分析知道，一点邻域内的运动由以下三项之和组成：\textcircled{1}各向同性的体积膨胀率$\nabla\cdot\boldsymbol{U}=\dot{\epsilon}$；\textcircled{2}体积不变的纯变形率$S_{ij}'$；\textcircled{3}准刚体角速度运动$\dfrac{\boldsymbol{\omega}}{2}$.从运动学角度来看，体膨胀率、纯变形率和准刚体角速度分别给出了速度场的信息.在刚体运动学中，给定刚体的角速度就可以确定由刚体转动产生的各点线速度.现在对流体运动提出一个类似的问题，给定流场的准刚体角速度场，它是否能唯一确定流体的速度场？如果它不能唯一确定速度场，那么在什么条件下才能唯一确定速度场呢？以下唯一性定理回答我们提出的问题.

\subsection*{1. 由速度场散度和旋度确定速度场的唯一性定理}

\noindent\textbf{定理2.4}\quad 已知域内速度场的散度和旋度以及边界上的法向速度，则唯一确定域内速度场.

按定理的前提，域内$D$中给定体积膨胀率$\theta(x_1,x_2,x_3)$和旋度$\boldsymbol{\omega}(x_1,x_2,x_3)$，即
\begin{gather*}
\nabla\cdot\boldsymbol{U}=\theta(x_1,x_2,x_3)\tag{2.46a}\\[2pt]
\nabla\times\boldsymbol{U}=\boldsymbol{\omega}(x_1,x_2,x_3)\tag{2.46b}
\end{gather*}
在边界$\Sigma$上给定$U_{bn}(x_1,x_2,x_3)$，即
\begin{equation*}
\boldsymbol{U}\cdot\boldsymbol{n}\,\big|_{\Sigma}=U_{bn}(x_1,x_2,x_3)
\tag{2.46c}
\end{equation*}
可以证明，在边界条件(2.46c)下方程(2.46a)和(2.46b)的解是唯一的.

\noindent\textbf{证明}\quad 设有两个速度场$\boldsymbol{U}_1$，$\boldsymbol{U}_2$同时满足域内方程和边界值.现在要证明$\boldsymbol{U}_1\equiv\boldsymbol{U}_2$.因为方程和边界条件都是线性的，所以$\boldsymbol{u}=\boldsymbol{U}_1-\boldsymbol{U}_2$满足以下方程和边界条件：
\begin{equation*}
\nabla\cdot\boldsymbol{u}=0,\qquad \nabla\times\boldsymbol{u}=\boldsymbol{0},\qquad \boldsymbol{u}\cdot\boldsymbol{n}\,\big|_{\Sigma}=0
\tag{2.47}
\end{equation*}
$\boldsymbol{u}$的旋度等于零，因此存在速度势$\boldsymbol{u}=\nabla\Phi$，代入散度公式得速度势满足拉普拉斯方程：
\[
\Delta\Phi=0
\]
及边界条件：$\boldsymbol{u}\cdot\boldsymbol{n}\,\big|_{\Sigma}=\dfrac{\partial\Phi}{\partial n}\Big|_{\Sigma}=0$.因此求解$\boldsymbol{u}=\boldsymbol{U}_1-\boldsymbol{U}_2$的边值问题(2.47)就是解满足第二类齐次边界条件的拉普拉斯方程，根据数学物理方程，它的解是常量，即$\Phi=\mathrm{const.}$，因而
\[
\boldsymbol{u}=\nabla\Phi=\boldsymbol{0}
\]
也就是$\boldsymbol{U}_1\equiv\boldsymbol{U}_2$，证毕.

上述唯一性定理说明，仅有域内速度场的旋度(或流体的准刚体角速度场)不能唯一确定流体速度场.按照唯一性定理要求，必须给定域内的速度散度、旋度以及边界上的法向速度才能唯一确定流体速度场.

由于散度方程、旋度方程(2.46a)，(2.46b)及边界条件(2.46c)都是线性的，可以将求解速度场的问题分解成三个部分求解.令
\begin{equation*}
\boldsymbol{U}=\boldsymbol{U}_e+\boldsymbol{U}_v+\boldsymbol{U}_p
\tag{2.48}
\end{equation*}
(1)$\boldsymbol{U}_e$是无旋有散度速度场的一个特解，即它满足
\begin{equation*}
\nabla\cdot\boldsymbol{U}_e=\theta,\qquad \nabla\times\boldsymbol{U}_e=\boldsymbol{0}
\tag{2.49}
\end{equation*}
(2)$\boldsymbol{U}_v$是有旋无散度速度场的一个特解，即它满足
\begin{equation*}
\nabla\cdot\boldsymbol{U}_v=0,\qquad \nabla\times\boldsymbol{U}_v=\boldsymbol{\Omega}
\tag{2.50}
\end{equation*}
(3)$\boldsymbol{U}_p$是无旋无散度流场，并满足物面不可穿透边界条件的解，即
\begin{gather*}
\nabla\cdot\boldsymbol{U}_p=0,\qquad \nabla\times\boldsymbol{U}_p=\boldsymbol{0}\tag{2.51a}\\[2pt]
\boldsymbol{U}_p\cdot\boldsymbol{n}\,\big|_{\Sigma}=U_{bn}-\boldsymbol{U}_e\cdot\boldsymbol{n}\,\big|_{\Sigma}-\boldsymbol{U}_v\cdot\boldsymbol{n}\,\big|_{\Sigma}\tag{2.51b}
\end{gather*}
利用式(2.49)、式(2.50)和式(2.51)，容易证明式$\boldsymbol{U}=\boldsymbol{U}_e+\boldsymbol{U}_v+\boldsymbol{U}_p$是满足方程(2.46a)和(2.46b)及边界条件(2.46c)的解，而唯一性定理证明该解是唯一的.下面我们给出以上第一和第二个问题的求解方法和结果.

\subsection*{2. 给定流场散度求速度场的特解}

由式(2.49)确定的流场只有散度，而无旋度，因此速度场有势：$\boldsymbol{U}_e=\nabla\Phi_e$，将它代入散度方程$\nabla\cdot\boldsymbol{U}_e=\theta$，速度势$\Phi_e$满足如下的泊松方程：
\[
\Delta\Phi_e=\theta(x_1,x_2,x_3)
\]
在数学物理方程中已给出，泊松方程的一个特解是
\begin{equation*}
\Phi_e(x_1,x_2,x_3)=-\frac{1}{4\pi}\iiint_D \frac{\theta(\xi,\eta,\zeta)}{R(x_1,x_2,x_3;\xi,\eta,\zeta)}\,\mathrm{d}\xi\,\mathrm{d}\eta\,\mathrm{d}\zeta
\tag{2.52a}
\end{equation*}
式中$(\xi,\eta,\zeta)$是积分域$D$内的积分变量，$(x_1,x_2,x_3)$是空间任意点的坐标，在积分号下$(x_1,x_2,x_3)$是常参量；$R=\big[(x_1-\xi)^{2}+(x_2-\eta)^{2}+(x_3-\zeta)^{2}\big]^{1/2}$是流场任意点$(x_1,x_2,x_3)$到$(\xi,\eta,\zeta)$之间的距离.由速度势$\Phi_e$计算$\boldsymbol{U}_e$的公式为
\[
\boldsymbol{U}_e=\nabla\Phi_e=\nabla\Big(-\frac{1}{4\pi}\iiint_D \frac{\theta(\xi,\eta,\zeta)}{R}\,\mathrm{d}\xi\,\mathrm{d}\eta\,\mathrm{d}\zeta\Big)=-\frac{1}{4\pi}\iiint_D \theta(\xi,\eta,\zeta)\,\nabla\frac{1}{R}\,\mathrm{d}\xi\,\mathrm{d}\eta\,\mathrm{d}\zeta
\]
式中$\nabla=\boldsymbol{e}_1\dfrac{\partial}{\partial x_1}+\boldsymbol{e}_2\dfrac{\partial}{\partial x_2}+\boldsymbol{e}_3\dfrac{\partial}{\partial x_3}$，它是对变量$\boldsymbol{x}=(x_1,x_2,x_3)$的向量导数，因此可移到积分号内.由向量求导公式，$\nabla\Big(\dfrac{1}{R}\Big)=-\dfrac{\boldsymbol{R}}{R^{3}}$，因而得
\begin{equation*}
\boldsymbol{U}_e=\frac{1}{4\pi}\iiint_D \theta(\xi,\eta,\zeta)\,\frac{\boldsymbol{R}}{R^{3}}\,\mathrm{d}\xi\,\mathrm{d}\eta\,\mathrm{d}\zeta
\tag{2.52b}
\end{equation*}
式中$\boldsymbol{R}=(x_1-\xi)\boldsymbol{e}_1+(x_2-\eta)\boldsymbol{e}_2+(x_3-\zeta)\boldsymbol{e}_3$.

\noindent\textbf{例2.10}\quad 设速度场散度$\theta(\xi,\eta,\zeta)$是$\delta$函数，即$\displaystyle\iiint_{D_0}\theta(\xi,\eta,\zeta)\,\mathrm{d}\xi\,\mathrm{d}\eta\,\mathrm{d}\zeta=1$，$D_0$是包含原点$(0,0,0)$的任意域；当$(\xi,\eta,\zeta)\neq(0,0,0)$时，$\theta(\xi,\eta,\zeta)=0$.求该速度场散度产生的速度场.

\noindent\textbf{解}\quad 将给定的速度场散度代入式(2.52)，积分式等于
\[
\iiint_D \theta(\xi,\eta,\zeta)\,\frac{\boldsymbol{R}}{R^{3}}\,\mathrm{d}\xi\,\mathrm{d}\eta\,\mathrm{d}\zeta=\Big(\frac{\boldsymbol{R}}{R^{3}}\Big)_{\!\xi=0,\,\eta=0,\,\zeta=0}=\frac{x_1\boldsymbol{e}_1+x_2\boldsymbol{e}_2+x_3\boldsymbol{e}_3}{(x_1^{2}+x_2^{2}+x_3^{2})^{3/2}}
\]
速度场为
\[
\boldsymbol{U}_e=\frac{1}{4\pi}\,\frac{\boldsymbol{R}}{R^{3}}
\]
该例题是不可压缩流场的一个典型解.除原点以外，该速度场处处散度为零，因此是不可压缩流场，但是在原点处有一个无限大的速度散度，也就是说在原点处有流体源源不断地流入流场，而流入的总流量是常量：$\displaystyle\iiint_{D_0}\theta(\xi,\eta,\zeta)\,\mathrm{d}\xi\,\mathrm{d}\eta\,\mathrm{d}\zeta=1$.在流体力学中这一典型解称为\textbf{点源}，后面还会详细分析和应用点源流场的解.

\subsection*{3. 由速度场旋度求速度场特解}

现在来求满足式(2.50)的速度场，该流场散度处处等于零，因此速度场可以表示为向量函数的旋度：
\[
\boldsymbol{U}_v=\nabla\times\boldsymbol{A}
\]
由于向量场旋度的散度等于零(参见附录)，以上给出的速度场满足$\nabla\cdot\boldsymbol{U}_v=0$，将$\boldsymbol{U}_v=\nabla\times\boldsymbol{A}$代入旋度方程，得
\[
\nabla\times\nabla\times\boldsymbol{A}=\boldsymbol{\omega}
\]
利用向量算符运算等式(见附录)：$\nabla\times\nabla\times\boldsymbol{A}=\nabla(\nabla\cdot\boldsymbol{A})-\nabla\cdot\nabla\boldsymbol{A}$，我们选择这样的特解，要求$\nabla\cdot\boldsymbol{A}=0$，于是$\boldsymbol{U}_v$的特解方程为
\[
\nabla\cdot\nabla\boldsymbol{A}=\nabla^{2}\boldsymbol{A}=-\boldsymbol{\omega}\,,\qquad \boldsymbol{U}_v=\nabla\times\boldsymbol{A}
\]
求出以上方程的特解后，再来验证是否满足$\nabla\cdot\boldsymbol{A}=0$，如果这一条件满足，特解就求出了.

导出的向量场$\boldsymbol{A}$满足向量型的泊松方程，也就是说，它的每一个分量都满足泊松方程.对每个旋度分量应用标量泊松方程解的公式(2.52a)，然后做向量和，就可得向量泊松方程的特解如下：
\[
\boldsymbol{A}(x_1,x_2,x_3)=\frac{1}{4\pi}\iiint_D \frac{\boldsymbol{\omega}(\xi,\eta,\zeta)}{R(x_1,x_2,x_3;\xi,\eta,\zeta)}\,\mathrm{d}\xi\,\mathrm{d}\eta\,\mathrm{d}\zeta
\]
对上式求旋度，得
\begin{align*}
\boldsymbol{U}_v=\nabla\times\boldsymbol{A}&=\frac{1}{4\pi}\,\nabla\times\iiint_D \frac{\boldsymbol{\omega}(\xi,\eta,\zeta)}{(x_1,x_2,x_3;\xi,\eta,\zeta)}\,\mathrm{d}\xi\,\mathrm{d}\eta\,\mathrm{d}\zeta\\
&=\frac{1}{4\pi}\iiint_D \nabla\times\frac{\boldsymbol{\omega}(\xi,\eta,\zeta)}{R}\,\mathrm{d}\xi\,\mathrm{d}\eta\,\mathrm{d}\zeta
\end{align*}
利用向量运算公式$\nabla\times(\psi\boldsymbol{\chi})=\psi\nabla\times\boldsymbol{\chi}-\boldsymbol{\chi}\times\nabla\psi$(见附录)，注意到向量算符$\nabla$只对$\boldsymbol{x}=(x_1,x_2,x_3)$求导数，故$\nabla_x\times\boldsymbol{\omega}(\xi,\eta,\zeta)=\boldsymbol{0}$，积分号下的向量运算等于：
\[
\nabla\times\frac{\boldsymbol{\omega}(\xi,\eta,\zeta)}{R}=-\boldsymbol{\omega}(\xi,\eta,\zeta)\times\nabla\frac{1}{R}=\frac{\boldsymbol{\omega}\times\boldsymbol{R}}{R^{3}}
\]
最后，由旋度场产生的速度场公式为
\begin{equation*}
\boldsymbol{U}_v=\frac{1}{4\pi}\iiint_D \frac{\boldsymbol{\omega}(\xi,\eta,\zeta)\times\boldsymbol{R}}{R^{3}}\,\mathrm{d}\xi\,\mathrm{d}\eta\,\mathrm{d}\zeta
\tag{2.53}
\end{equation*}
上式是著名的毕奥-萨伐尔(Biot-Savart)公式，该式和由电流诱导磁场的公式相仿，因此，式(2.53)常称作涡诱导速度公式.最后，需要验证已求出的$\boldsymbol{A}$是否满足给定的条件$\nabla\times\boldsymbol{A}=\boldsymbol{0}$.

令$\nabla'=\boldsymbol{e}_1\dfrac{\partial}{\partial\xi}+\boldsymbol{e}_2\dfrac{\partial}{\partial\eta}+\boldsymbol{e}_3\dfrac{\partial}{\partial\zeta}$，则对于$R=\sqrt{(x_1-\xi)^{2}+(x_2-\eta)^{2}+(x_3-\zeta)^{2}}$应有$\nabla\Big(\dfrac{1}{R}\Big)=-\nabla'\Big(\dfrac{1}{R}\Big)$.将$\nabla\cdot\boldsymbol{A}$展开，并在运算过程中应用上述关系，就有
\begin{align*}
\nabla\cdot\boldsymbol{A}&=\frac{1}{4\pi}\,\nabla\cdot\iiint_D \frac{\boldsymbol{\omega}'}{R}\,\mathrm{d}\xi\,\mathrm{d}\eta\,\mathrm{d}\zeta=\frac{1}{4\pi}\iiint_D \boldsymbol{\omega}'\cdot\nabla\frac{1}{R}\,\mathrm{d}\xi\,\mathrm{d}\eta\,\mathrm{d}\zeta=-\frac{1}{4\pi}\iiint_D \boldsymbol{\omega}'\cdot\nabla'\frac{1}{R}\,\mathrm{d}\xi\,\mathrm{d}\eta\,\mathrm{d}\zeta\\
&=-\frac{1}{4\pi}\iiint_D \nabla'\cdot\Big(\frac{\boldsymbol{\omega}'}{R}\Big)\mathrm{d}\xi\,\mathrm{d}\eta\,\mathrm{d}\zeta+\frac{1}{4\pi}\iiint_D \frac{1}{R}\,\nabla'\cdot\boldsymbol{\omega}'\,\mathrm{d}\xi\,\mathrm{d}\eta\,\mathrm{d}\zeta
\end{align*}
式中$\boldsymbol{\omega}'$是$\boldsymbol{\omega}(\xi,\eta,\zeta)$的简写.因为旋度的散度等于零，上面公式的最后一项等于零.第一个积分式可用高斯公式化成面积分：
\[
\iiint_D \nabla'\cdot\frac{\boldsymbol{\omega}'}{R}\,\mathrm{d}\xi\,\mathrm{d}\eta\,\mathrm{d}\zeta=\iint_\Sigma \frac{\boldsymbol{\omega}\cdot\boldsymbol{n}}{R}\,\mathrm{d}A
\]
于是
\[
\nabla\cdot\boldsymbol{A}=-\frac{1}{4\pi}\iint_\Sigma \frac{\boldsymbol{\omega}\cdot\boldsymbol{n}}{R'}\,\mathrm{d}A
\]
由上式可见，如果在求解域的边界上，处处满足$\boldsymbol{\omega}\cdot\boldsymbol{n}=0$，则在域内满足$\nabla\cdot\boldsymbol{A}=0$.以下两种情况可以满足上述条件.

(1)求解域边界是涡面，因为涡面上$\boldsymbol{\omega}\cdot\boldsymbol{n}=0$.

例如，无界域中孤立涡管周围的速度场，可用式(2.53)计算.

(2)静止固壁的粘附边界.因为静止固壁上速度等于零，从而固壁上任意封闭周线的速度环量也等于零：$\displaystyle\oint\boldsymbol{U}\cdot\mathrm{d}\boldsymbol{r}=0$，再利用斯托克斯公式，在固壁上任意封闭曲线所张的表面$A$上有：$\displaystyle\iint_A \boldsymbol{\omega}\cdot\boldsymbol{n}\,\mathrm{d}A=\oint\boldsymbol{U}\cdot\mathrm{d}\boldsymbol{r}=0$.积分周线是任取的，因此在固壁表面处处满足$\boldsymbol{\omega}\cdot\boldsymbol{n}=0$.

所以，在静止固壁包围的有界域内，上述的涡诱导公式(2.52)也可以适用.

\subsection*{4. 满足固壁法向速度条件的无旋、无散度速度场的解}

最后，我们简单地说明满足固壁法向速度条件的无旋、无散度速度场的解法，该问题的方程和边界条件为前面导出的式(2.51a)和式(2.51b)：
\begin{gather*}
\nabla\cdot\boldsymbol{U}_p=0,\qquad \nabla\times\boldsymbol{U}_p=\boldsymbol{0}\\[2pt]
\boldsymbol{U}_p\cdot\boldsymbol{n}\,\big|_{\Sigma}=U_{bn}-\boldsymbol{U}_e\cdot\boldsymbol{n}\,\big|_{\Sigma}-\boldsymbol{U}_v\cdot\boldsymbol{n}\,\big|_{\Sigma}
\end{gather*}
由于无旋速度场存在速度势$\Phi$，$\boldsymbol{U}_p=\nabla\Phi$，该速度场又要满足散度等于零，因此速度势必须满足拉普拉斯方程：
\[
\Delta\Phi=0
\]
在所给边界条件下(2.51b)，该问题属于拉普拉斯方程的第2类边界条件问题，在数理方程中有详细讨论，本书将在第5章中作较详细的介绍.

\subsection*{5. 由涡量场求速度场的算例}

(1)直涡线的诱导速度场

设在无界流场中有一无限长的细直涡管，涡管强度为$\Gamma$，求该涡管周围的诱导速度，参见图2.13(a).

利用毕奥-萨伐尔公式：
\[
\boldsymbol{U}_v=\frac{1}{4\pi}\iiint_D \frac{\boldsymbol{\omega}(\xi,\eta,\zeta)\times\boldsymbol{R}}{R^{3}}\,\mathrm{d}\xi\,\mathrm{d}\eta\,\mathrm{d}\zeta
\]
设直角坐标的$z$轴和涡线一致，则$\displaystyle\iiint_u \boldsymbol{\omega}\,\mathrm{d}\xi\,\mathrm{d}\eta=\Gamma\boldsymbol{e}_3$，$\boldsymbol{R}=x_1\boldsymbol{e}_1+x_2\boldsymbol{e}_2+(x_3-\zeta)\boldsymbol{e}_3$，因此涡诱导速度为
\begin{align*}
\boldsymbol{U}_v&=\frac{1}{4\pi}\int_{-\infty}^{+\infty} \frac{\Gamma\boldsymbol{e}_3\times\boldsymbol{R}}{R^{3}}\,\mathrm{d}\zeta=\frac{\Gamma}{4\pi}\int_{-\infty}^{+\infty} \frac{\boldsymbol{e}_3\times\boldsymbol{R}}{R^{3}}\,\mathrm{d}\zeta\\
&=\frac{\Gamma}{4\pi}\int_{-\infty}^{+\infty} \frac{x_1\boldsymbol{e}_2-x_2\boldsymbol{e}_1}{\big[x_1^{2}+x_2^{2}+(x_3-\zeta)^{2}\big]^{3/2}}\,\mathrm{d}\zeta
\end{align*}
积分上式后，得
\[
\boldsymbol{U}_v=\frac{\Gamma}{2\pi}\bigg(\frac{-x_2\boldsymbol{e}_1+x_1\boldsymbol{e}_2}{x_1^{2}+x_2^{2}}\bigg)
\]
以上结果表明：无限长直涡线的诱导速度场是垂直于涡线的平面流场，流线是以涡线为轴心的圆周线(参见图2.13(b)).

(2)圆周涡线的诱导速度

设有涡量强度为$\Gamma$的平面圆涡线，半径为$R$，求圆心处的涡诱导速度.应用毕奥-萨伐尔公式，类似直涡线的计算，将体积分化为线积分，得
\[
\boldsymbol{U}(0,0,0)=\frac{\Gamma}{4\pi}\oint_c \frac{\mathrm{d}\boldsymbol{l}\times\boldsymbol{R}}{R^{3}}
\]
在圆涡线上：$\mathrm{d}\boldsymbol{l}=\boldsymbol{e}_\theta R\,\mathrm{d}\theta$，$\boldsymbol{R}=-R\boldsymbol{e}_r$，代入积分式，得
\[
\boldsymbol{U}(0,0,0)=\frac{\Gamma\boldsymbol{e}_z}{4\pi}\int_0^{2\pi}\frac{\mathrm{d}\theta}{R}=\frac{\Gamma}{2R}\boldsymbol{e}_z
\]
于是得到圆涡线的圆心处的诱导速度垂直于涡线平面(参见图2.13(b)).除圆心以外，平面圆周涡线诱导速度场的积分式不能用简单解析函数表示.

(3)空间曲线涡附近的诱导速度

在非均匀流场中，常见曲涡管.作为一种近似，研究一般空间曲线涡的诱导速度.空间曲线上任意一点的几何特性，可以用它的切线和法平面(垂直于切线)来描述，在法平面上垂直于切线的有互相垂直两条法线，分别称为主法线和次法线.空间曲线涡与直线涡一样在涡线的法平面内有一围绕涡线的渐近圆周运动，除此以外曲涡线还诱导次法线方向的速度，因此曲线涡不可能自身固定在流场中而要移动和变形，下面应用毕奥-萨伐尔公式对一般曲线涡某点邻域的速度场加以分析.

设曲线涡上一点$O$处的切线方向为$\boldsymbol{t}$，主法线方向为$\boldsymbol{n}$，次法线方向为$\boldsymbol{b}$.在$O$点垂直于曲线的法平面上任一点的位置向量为(见图2.14(b))：
\[
\boldsymbol{x}=x_2\boldsymbol{n}+x_3\boldsymbol{b}=\sigma\cos\phi\,\boldsymbol{n}+\sigma\sin\phi\,\boldsymbol{b}
\]
现在考察在$O$点邻域$((−L,L)$和$\sqrt{(x_2^{2}+x_3^{2})}=\sigma\to 0)$曲线涡诱导流场的渐近性质.$O$点附近弧段内涡线上的位置向量为
\[
\boldsymbol{\zeta}=l\,\boldsymbol{t}+\frac{1}{2}cl^{2}\boldsymbol{n}
\]
此处$c$为$O$点的曲率，$l$为切线方向$\boldsymbol{t}$的坐标量，因而$O$点附近，$\mathrm{d}\boldsymbol{x}=(\boldsymbol{t}+cl\,\boldsymbol{n})\,\mathrm{d}l$，从而
\[
\frac{\boldsymbol{R}\times\mathrm{d}\boldsymbol{x}}{R^{3}}=\frac{(\boldsymbol{x}-\boldsymbol{\zeta})\times\mathrm{d}\boldsymbol{x}}{|\boldsymbol{x}-\boldsymbol{\zeta}|^{3}}=\frac{-x_3cl\,\boldsymbol{t}+x_3\boldsymbol{n}-\Big(x_2+\dfrac{1}{2}cl^{2}\Big)\boldsymbol{b}}{\Big[x_3^{2}+x_2^{2}+l^{2}(1-x_2c)+\dfrac{1}{4}c^{2}l^{4}\Big]^{3/2}}\,\delta l
\]
在微元长度$(-L,L)$上积分上式，分子上$-x_3cl$的项关于$l$反对称，因此积分值等于零，即涡线没有切向的诱导速度.于是，涡线对$O$点$(0,\sigma\cos\phi,\sigma\sin\phi)$的诱导速度为
\[
\frac{1}{4\pi}\int_{-L/\sigma}^{L/\sigma} \frac{(\boldsymbol{b}\cos\phi-\boldsymbol{n}\sin\phi)\,\sigma^{-1}+\dfrac{1}{2}cm^{2}\boldsymbol{b}}{\Big[1+m^{2}(1-c\sigma\cos\phi)+\dfrac{1}{4}c^{2}\sigma^{2}m^{4}\Big]^{3/2}}\,\mathrm{d}m
\]
此处：$m=\dfrac{l}{\sigma}$，$\cos\phi=\dfrac{x_2}{\sigma}$，$\sin\phi=\dfrac{x_3}{\sigma}$，以上积分式(可应用积分公式表)在$\sigma\to 0$时的渐近值为
\begin{equation*}
\boldsymbol{U}_{\sigma\to 0}=\frac{\Gamma}{2\pi\sigma}(\boldsymbol{b}\cos\phi-\boldsymbol{n}\sin\phi)+\frac{\Gamma c}{2\pi}\,\boldsymbol{b}\ln\frac{L}{\sigma}+O(\sigma^{0})
\tag{2.54}
\end{equation*}
上式右边第一部分与直线涡周围产生的流体运动相同，流体质点围绕直涡管作圆周运动，并当$\sigma\to 0$时有奇性；第二部分的诱导速度在次法线方向，并具有对数奇性.也就是说曲线涡本身在次法线方向诱导很大的移动速度.以上分析表明，曲线涡本身不能在流场中保持静止，也不能保持它的几何形状，它在运动过程不断变形.

对于理想的线涡，在涡线上的诱导速度是无限大.实际流体有粘性，曲涡线是有一定半径的曲涡管，涡管内涡量由零逐渐增加，因此不存在速度奇性.如涡管截面很小，则曲线涡诱导速度公式(2.54)是很好近似.

对于平面圆周线涡环来说，涡线各点的曲率及次法线方向不变，$c=\dfrac{1}{R}$，$\boldsymbol{b}=\boldsymbol{i}$，因此整个涡环作平移运动，但它的移动速度远远大于涡环中心的速度.如果把坐标系固定在涡环上，或者说，在随涡环一起移动的坐标架中观察绕涡环的流动，将得到如图2.15所示的流线谱.

(4)片涡

如果有旋流体质点集中在一个面上，涡量必沿该面的切线方向，这种理想的旋涡称为片涡(也称涡片).

片涡可以看作为厚度极薄的涡层(图2.16(a))，片涡的曲面可以用双参数方程表示：
\begin{equation*}
\boldsymbol{x}=\boldsymbol{g}(l,m,t)
\tag{2.55}
\end{equation*}
片涡的涡量用它面上的涡量强度表示：
\begin{equation*}
\boldsymbol{\gamma}=\lim_{\substack{\varepsilon\to 0\\ \delta\to\infty}}\int_0^{\varepsilon}\boldsymbol{\omega}\,\mathrm{d}n
\tag{2.56}
\end{equation*}
式中$\boldsymbol{\gamma}$称为片涡的面强度，$\boldsymbol{\omega}$是片涡内的涡量分布，$\varepsilon$是片涡的厚度.

片涡周围的诱导速度也可用毕奥-萨伐尔公式计算如下：
\begin{equation*}
\boldsymbol{U}=-\frac{1}{4\pi}\int_V \frac{\boldsymbol{R}\times\boldsymbol{\omega}\,\mathrm{d}n\,\mathrm{d}A}{R^{3}}=-\frac{1}{4\pi}\int_A \frac{\boldsymbol{R}\times\boldsymbol{\gamma}\,\mathrm{d}A}{R^{3}}
\tag{2.57}
\end{equation*}
片涡是一种切向速度的间断面.为了说明这一点，可以在片涡上任选一点$O$，并在该点作法向量$\boldsymbol{n}$和切平面$T$(图2.16(b)).在切平面上以$O$为圆心作半径为$\eta$的圆，并在法线方向取$\delta$为高作一圆柱体，并考察柱顶点$M$处的速度$\boldsymbol{U}_M$.顶点$M$处$\boldsymbol{x}=\delta\boldsymbol{n}$，涡面上$\boldsymbol{\zeta}=\boldsymbol{r}$，故$\boldsymbol{R}=\delta\boldsymbol{n}-\boldsymbol{r}$，$R=\sqrt{\delta^{2}+r^{2}}$，代入式(2.57)，得
\begin{align*}
\boldsymbol{U}_M&=-\frac{1}{4\pi}\int_0^{\eta}\!\!\int_0^{2\pi} \frac{\delta\boldsymbol{n}\times\boldsymbol{\gamma}\,r\,\mathrm{d}r\,\mathrm{d}\theta}{(\delta^{2}+r^{2})^{3/2}}+\frac{1}{4\pi}\int_0^{\eta}\!\!\int_0^{2\pi} \frac{\boldsymbol{r}\times\boldsymbol{\gamma}\,r\,\mathrm{d}r\,\mathrm{d}\theta}{(\delta^{2}+r^{2})^{3/2}}\\
&=-\frac{1}{4\pi}\int_0^{\eta}\!\!\int_0^{2\pi} \frac{\delta\boldsymbol{n}\times\boldsymbol{\gamma}\,r\,\mathrm{d}r\,\mathrm{d}\theta}{(\delta^{2}+r^{2})^{3/2}}+\frac{1}{4\pi}\int_0^{\eta}\!\!\int_0^{2\pi} \frac{|\boldsymbol{r}|\times|\boldsymbol{\gamma}|\,\sin\theta\,r\,\mathrm{d}r\,\mathrm{d}\theta}{(\delta^{2}+r^{2})^{3/2}}\,\boldsymbol{n}\\
&=\frac{-\boldsymbol{n}\times\boldsymbol{\gamma}}{2}\int_0^{\eta} \frac{\delta r\,\mathrm{d}r}{(\delta^{2}+r^{2})^{3/2}}=\frac{-\boldsymbol{n}\times\boldsymbol{\gamma}}{4}\int_0^{\eta} \frac{\mathrm{d}\big(\dfrac{r}{\delta}\big)^{2}}{\Big(1+\dfrac{r^{2}}{\delta^{2}}\Big)^{3/2}}\\
&=\frac{1}{2}\,\boldsymbol{n}\times\boldsymbol{\gamma}\,\Big(1+\frac{r^{2}}{\delta^{2}}\Big)^{1/2}\bigg|_0^{\eta}=-\frac{1}{2}\,\boldsymbol{n}\times\boldsymbol{\gamma}\,\Big[1-\Big(1+\frac{\eta^{2}}{\delta^{2}}\Big)^{-1/2}\Big]
\end{align*}
注意积分式中$\boldsymbol{r}$是切平面上的径向，$\eta$是切平面上微圆的半径.令$\eta$和$\delta$同时趋于零，但是$\delta/\eta\to 0$，得在涡片正法线方向一面的诱导速度
\begin{equation*}
\boldsymbol{U}_{+}=\lim_{\delta/\eta\to 0}\boldsymbol{U}_M=\frac{1}{2}\,\boldsymbol{\gamma}\times\boldsymbol{n}
\tag{2.58}
\end{equation*}
同理涡片另一侧的诱导速度为
\begin{equation*}
\boldsymbol{U}_{-}=-\frac{1}{2}\,\boldsymbol{\gamma}\times\boldsymbol{n}
\tag{2.59}
\end{equation*}
由于$\boldsymbol{\gamma}$与涡片相切，$\boldsymbol{n}$与涡片垂直，故$\boldsymbol{U}_{+}$，$\boldsymbol{U}_{-}$与涡片相切，在涡片上切向速度的间断量
\begin{equation*}
[\boldsymbol{U}]=\boldsymbol{U}_{+}-\boldsymbol{U}_{-}=\boldsymbol{\gamma}\times\boldsymbol{n}
\tag{2.60}
\end{equation*}
取式(2.60)的绝对值，得：$|\boldsymbol{U}_{+}-\boldsymbol{U}_{-}|=\gamma$，也就是说：\textbf{片涡强度等于涡片两侧的切向速度差}.将式(2.60)两侧叉乘$\boldsymbol{n}$，得到用切向速度间断量来表示的片涡的面涡强度
\begin{equation*}
\boldsymbol{\gamma}=-[\boldsymbol{U}]\times\boldsymbol{n}
\tag{2.61}
\end{equation*}

\subsection*{6. Föppl定理和应用}

涡量场中一些固有的运动学特性对于研究流体运动是很有用的.Föppl定理对于研究绕流物体的涡量场有很大帮助.

(1)Föppl定理：\textbf{在平移刚体运动的无界流场中，全场总涡量等于零.}

证明如下：根据涡量的定义，很容易导出以下等式：
\begin{equation*}
\omega_i=\frac{\partial}{\partial x_j}\,x_i\omega_j
\tag{2.62}
\end{equation*}
上式右边：$\dfrac{\partial}{\partial x_j}x_i\omega_j=\omega_j\dfrac{\partial x_i}{\partial x_j}+x_i\dfrac{\partial\omega_j}{\partial x_j}$，由于涡量场散度等于零，同时$\partial x_i/\partial x_j=\delta_{ij}$，等式右边简化为
\[
\frac{\partial}{\partial x_j}x_i\omega_j=\omega_j\delta_{ij}=\omega_i
\]
式(2.62)证毕.将式(2.62)在全场积分，并用高斯公式将体积分化为面积分得
\begin{align*}
\int_D \omega_i\,\mathrm{d}V&=\int_D \Big(\frac{\partial}{\partial x_j}x_i\omega_j\Big)\mathrm{d}V=\oint_{\partial D}(x_i\omega_j n_j)\,\mathrm{d}A\\
&=\oint_{A_w}(x_i\omega_j n_j)\,\mathrm{d}A+\oint_{A_\infty}(x_i\omega_j n_j)\,\mathrm{d}A
\tag{2.63}
\end{align*}
$\partial D$是全场积分的边界，第1个积分在绕流固壁上求和，因为平移运动壁面的法向涡量等于零(见前面3.(2)小节)；第2个积分包含绕流物体的无穷大球体中求和，因为无穷远处$\omega\sim O(r^{-4})$也等于零，于是有
\begin{equation*}
\int_D \omega_i\,\mathrm{d}V=0
\tag{2.64}
\end{equation*}
式(2.64)的意义是，由刚体平移运动产生的无界流场中，涡量的总量始终等于零，不随时间变化.如果说在某一区域中产生了正旋涡，则必在流场的另一处有总量相等的负旋涡.这是涡量运动学的性质，与流体性质无关，在理想流体、粘性流体、不可压缩流体和可压缩流体中都成立.

例如，流体绕圆柱或圆球流动的尾迹中常常出现一系列旋涡，而这些旋涡必定成对出现，例如著名的卡门涡街(见图2.17).另外，在绕机翼的有攻角流动中，机翼产生升力，用理想流体理论可以证明(详见第5章)围绕翼型的周线上必有环量，根据Föppl定理，在流场中必产生旋涡，其环量和绕物体环量方向相反，在空气动力学中，与绕物体环量方向相反的旋涡称为启动涡(见图2.18).

关于涡量的其他运动学性质，可参见涡动力学的专著.
\section*{习题}\addcontentsline{toc}{section}{习题}

\noindent\textbf{2.1}\quad 已知拉格朗日速度分布$u=-\beta(a\sin\beta t+b\cos\beta t)$，$v=\beta(a\cos\beta t-b\sin\beta t)$，$w=\alpha$.如$t=0$时$x=a$，$y=b$，$z=c$，试用欧拉变量表示上述流场速度分布，并求流场加速度分布，式中$\alpha,\beta,a,b,c$为常数.

\medskip
\noindent\textbf{2.2}\quad 给定速度场$u=x+y$，$v=x-y$，$w=0$，且令$t=0$时$x=a$，$y=b$，$z=c$，求质点空间分布.

\medskip
\noindent\textbf{2.3}\quad 给定柱坐标中速度场的欧拉表达式：$U_r=\dfrac{t}{r}$，$U_\theta=\dfrac{1}{r}$，$U_z=0$，且令$t=0$时$r=a$，$\theta=b$，$z=c$，求质点分布及流场加速度的欧拉表达式.

\medskip
\noindent\textbf{2.4}\quad 已知平面速度场$u=x+t$，$v=-y+t$，并令$t=0$时$x=a$，$y=b$，求：

(1)流线方程及$t=0$时过$(-1,-1)$点的流线；

(2)迹线方程及$t=0$时过$(-1,-1)$点的迹线；

(3)用拉格朗日变量表示速度分布.

\medskip
\noindent\textbf{2.5}\quad 已知平面速度$u=1+2t$，$v=3+4t$，求：

(1)流线方程；

(2)$t=0$时，经过$(0,0)$，$(0,1)$，$(0,-1)$点的三条流线方程；

(3)$t=0$时，位置在$(0,0)$点的流体质点的迹线方程.

\medskip
\noindent\textbf{2.6}\quad 给定速度场$u=-ky$，$v=kx$，$w=w_0$，求通过$x=a$，$y=b$，$z=c$点的流线，式中$k,w_0$是常数.

\medskip
\noindent\textbf{2.7}\quad 推导下述平面流动的流线方程：$u=u_0$，$v=v_0\cos(kx-at)$，式中$v_0,k$和$a$均为常数.求$t=0$时通过$x=0,y=0$点的流线和迹线方程.当$k,a\to 0$时，比较这两条流线和迹线.

\medskip
\noindent\textbf{2.8}\quad 如果流体质点在圆锥面上运动，其速度为$U_R=A$，$U_\theta=0$，$U_\phi=B\cos\phi$，其中$A,B$都是常数，求$t=0$时通过$R=a,\theta=b,\phi=c$点的迹线和流线，并用拉格朗日变量表示此速度场，$R,\theta,\phi$分别是球坐标的向径、纬度角和经度角.

\medskip
\noindent\textbf{2.9}\quad 设$\boldsymbol{U}\neq\boldsymbol{0}$，说明以下三种导数
\[
\frac{D\boldsymbol{U}}{Dt}=\boldsymbol{0},\qquad \frac{\partial\boldsymbol{U}}{\partial t}=\boldsymbol{0},\qquad (\boldsymbol{U}\cdot\nabla)\boldsymbol{U}=0
\]
的物理意义.

\medskip
\noindent\textbf{2.10}\quad 判断下列两个问题，并说明理由.

(1)下述用拉格朗日变量表示的流体运动：
\[
x=a+bt+\frac{1}{2}ct^{2},\qquad y=b+ct,\qquad z=c
\]
是否为定常流动？

(2)在下述用平面极坐标系$(r,\theta)$表示的速度场：
\[
U_r=U_\infty\Big(1-\frac{a^{2}}{r^{2}}\Big)\cos\theta,\qquad U_\theta=-U_\infty\Big(1+\frac{a^{2}}{r^{2}}\Big)\sin\theta
\]
式中$U_\infty$和$a$都是常数，$r=a$是否为流线？

\medskip
\noindent\textbf{2.11}\quad 给定拉格朗日表达式的运动规律：
\[
x=a\cos\Big(\frac{t}{\sqrt{a^{2}+b^{2}}}\Big)-b\sin\Big(\frac{t}{\sqrt{a^{2}+b^{2}}}\Big),\qquad
y=b\cos\Big(\frac{t}{\sqrt{a^{2}+b^{2}}}\Big)+a\sin\Big(\frac{t}{\sqrt{a^{2}+b^{2}}}\Big)
\]
(1)求$t=0$时，通过点$(x=1,y=\sqrt{3})$的迹线方程和流线方程；

(2)它是否是定常流动？

(3)它是否是有旋流动？

\medskip
\noindent\textbf{2.12}\quad 若已知温度场$T=\dfrac{A}{x^{2}+y^{2}+z^{2}}\,t^{2}$和速度分布$u=xt$，$v=yt$，$w=zt$，现有一流体质点在该流场中运动，质点在$t=0$时的位置为$x=a$，$y=b$，$z=c$，试求该流体质点的温度随时间的变化关系，式中$A$为常数.

\medskip
\noindent\textbf{2.13}\quad 在球坐标和柱坐标中，试用欧拉变量表示常比热容完全气体的质点沿轨迹熵值不变的关系式.已知熵$s=c_v\ln\dfrac{p}{\rho^{\gamma}}+\mathrm{const.}$

\medskip
\noindent\textbf{2.14}\quad 给定拉格朗日位移表达式：$x=ae^{-\frac{2t}{K}}$，$y=b\Big(1+\dfrac{t}{K}\Big)^{2}$，$z=ce^{\frac{2t}{K}}\Big(1+\dfrac{t}{K}\Big)^{-2}$，式中$K$为常数$(K\neq 0)$，$a,b,c$分别为$t=0$时$x,y,z$的坐标值.请判别：(1)该流场是否为定常流场？(2)是否不可压流场？(3)是否有旋流场？

\medskip
\noindent\textbf{2.15}\quad 给定拉格朗日位移表达式：$x=ae^{-\frac{2t}{K}}$，$y=be^{\frac{t}{K}}$，$z=ce^{\frac{t}{K}}$，式中$K$为常数，$a,b,c$分别为$t=0$时$x,y,z$的坐标值.请判别：(1)该流场是否为定常流场？(2)是否不可压流场？(3)是否无旋流场？

\medskip
\noindent\textbf{2.16}\quad 给定柱坐标中速度场的欧拉表达式：
\[
U_r=U_\infty\Big(1-\frac{a^{2}}{r^{2}}\Big)\cos\theta,\qquad U_\varphi=-U_\infty\Big(1+\frac{a^{2}}{r^{2}}\Big)\sin\theta,\qquad U_z=0
\]
式中$U_\infty$和$a$为常数.

(1)试证$r=a$是流线；

(2)求$t=0$时位于$\theta=\pi$，$r=b\,(b>a)$的流体质点到达$\theta=\pi$，$r=a$需要的时间；

(3)求$t=0$时位于$\theta=\dfrac{\pi}{2}$，$r=a$的流体质点到达$\theta=0$，$r=a$所需要的时间.

\medskip
\noindent\textbf{2.17}\quad 求出下述问题的变形速率分量、涡量和体膨胀速率.并说明运动是否有旋？流体是否不可压缩？

(1)直角坐标系$(x,y,z)$内
\[
u=U(h^{2}-y^{2}),\qquad v=w=0
\]
式中$h,U$均为常量.

(2)柱坐标系$(r,\theta,z)$内
\[
U_z=U\bigg[\frac{a^{2}\ln(r/b)-b^{2}\ln(r/a)}{\ln(a/b)}-r^{2}\bigg],\qquad U_\theta=U_r=0
\]
式中$U,a,b$均为常量.

(3)球坐标系$(R,\theta,\phi)$内
\[
U_R=U_\infty\Big[1-\Big(\frac{a}{R}\Big)^{3}\Big]\cos\theta,\qquad
U_\theta=-U_\infty\Big[1+\frac{1}{2}\Big(\frac{a}{R}\Big)^{3}\Big]\sin\theta,\qquad U_\phi=0
\]
式中$U_\infty,a$均为常量.

\medskip
\noindent\textbf{2.18}\quad 给定不定常流动速度场
\begin{gather*}
U=U_0\bigg[1-\frac{2}{\sqrt{\pi}}\int_0^{\eta}\mathrm{e}^{-\eta^{2}}\mathrm{d}\eta\,\bigg],\qquad \eta=\frac{y}{2\sqrt{\nu t}}\\[2pt]
V=0
\end{gather*}
式中$\nu$为运动粘性系数(常数)，$U_0$为常数，求涡量场.

\medskip
\noindent\textbf{2.19}\quad 给定速度场$u=ax$，$v=ay$，$w=-2az$，式中$a$为常数.求：

(1)线变形速率分量、角变形速率分量和体积膨胀率；

(2)该流场是否是无旋流场？若无旋，写出它的速度势函数.

\medskip
\noindent\textbf{2.20}\quad 在欧拉描述的速度场$\boldsymbol{U}=\boldsymbol{U}(x,y,z,t)$中，请证明：

(1)变形率张量为零的有旋运动必为
\[
\boldsymbol{U}=\boldsymbol{U}_0+\boldsymbol{\omega}\times\boldsymbol{r}
\]
其中$\boldsymbol{U}_0,\boldsymbol{\omega}$与向量$\boldsymbol{r}(x,y,z)$无关；

(2)除了均匀流场(即$\boldsymbol{U}=$常向量)外，一切无旋流动的变形率张量都不等于零.

\medskip
\noindent\textbf{2.21}\quad 已知速度场为$u=y+2z$，$v=z+2x$，$w=x+2y$，

(1)求涡量及涡线方程；

(2)求在$x+y+z=1$平面上通过横截面积$\mathrm{d}A=0.0001\,\mathrm{m^2}$的涡管强度；

(3)求在$z=0$平面上通过横截面积$\mathrm{d}A=0.0001\,\mathrm{m^2}$的涡通量.

\medskip
\noindent\textbf{2.22}\quad 给定一个函数$\Phi=U_\infty\Big(r\cos\theta+\dfrac{1}{2}\ln r\Big)$，$r,\theta$是平面极坐标，问：

(1)这个函数能否作为一种二维不可压流动的速度势函数？

(2)如能，求速度场；

(3)计算$r=\sqrt{2}$，$\theta=\dfrac{\pi}{4}$点沿流线方向的加速度分量.

\medskip
\noindent\textbf{2.23}\quad 对二维不可压缩流体运动$(w=0)$，试证明

(1)如运动为无旋，则必满足$\nabla^{2}u=0$，$\nabla^{2}v=0$；

(2)满足$\nabla^{2}u=0$，$\nabla^{2}v=0$的流动不一定无旋.

\medskip
\noindent\textbf{2.24}\quad 给定柱坐标内流场速度势$\Phi=U_\infty(t)\cos\theta\,(r+a^{2}/r)$，求流线方程，并求通过$r=a,\theta=\pi/2$点的流线.

\medskip
\noindent\textbf{2.25}\quad 原静止不可压无界流场中，在位于$z=0$平面上放置强度为$\Gamma$的$\Pi$形涡线(如图2.19所示)，求其速度场.

\medskip
\noindent\textbf{2.26}\quad 原静止无界无旋流场中，给定$\nabla\cdot\boldsymbol{U}=q$，当$a\leqslant r\leqslant b$时，$q=\mathrm{const.}$；当$r>b$或$r<a$时，$q=0$.式中$a,b$为常数，$r$为柱坐标系中的径向坐标，求速度场.

提示：在柱坐标中求解，并注意$r=b$的圆周上速度连续.

\medskip
\noindent\textbf{2.27}\quad 相距为$h$的两个平面圆周线涡，圆周半径相同，均为$a$，线涡强度也相同，均为$\Gamma$，求线涡圆心处的诱导速度.

\section*{思考题}\addcontentsline{toc}{section}{思考题}

\noindent\textbf{S2.1}\quad 为什么质点位置变量$\boldsymbol{x}$是初始时刻位置变量$\boldsymbol{A}$一一对应的连续函数？

\medskip
\noindent\textbf{S2.2}\quad 从消防喷嘴喷出的水柱表面是流体面(或称“流体质面”)还是流面？流面是否可能同时是流体面？

\medskip
\noindent\textbf{S2.3}\quad 哪种运动中流体微团只有转动而无变形？哪种流体运动只有变形而无转动？是否存在既无变形又无旋转的流场？

\medskip
\noindent\textbf{S2.4}\quad 是否存在涡线和流线相互垂直的流动，试举例说明；是否存在流线和涡线平行的流动，试举例说明.

\medskip
\noindent\textbf{S2.5}\quad 在流场中孤立的曲线(面)涡管(涡管外涡量等于零)能否运动？孤立的直线涡管能否运动？为什么？在流场中孤立的曲面涡片(涡片外涡量等于零)能否运动？为什么？
